\documentclass[11pt,a4paper]{article}
\usepackage[margin=23mm]{geometry}
\usepackage[T1]{fontenc}
\usepackage{lmodern,microtype,amsmath,amssymb,amsthm,mathtools,enumitem,xcolor,booktabs}
\usepackage[colorlinks=true,linkcolor=blue!45!black,citecolor=blue!45!black,urlcolor=blue!45!black]{hyperref}
\newtheorem{theorem}{Theorem}[section]
\newtheorem{lemma}[theorem]{Lemma}
\newtheorem{proposition}[theorem]{Proposition}
\newtheorem{corollary}[theorem]{Corollary}
\theoremstyle{remark}\newtheorem{remark}[theorem]{Remark}
\setlength{\parskip}{3pt}
\setlist{nosep}
\allowdisplaybreaks[1]
\title{Dilation structure, uniform profile gaps,\\and their failure under small dispersion changes}
\author{Research continuation: analytic arguments and internal independent review}
\date{8 October 2026}
\begin{document}
\maketitle
\begin{abstract}
For a mass-critical focusing equation with a real-even positive elliptic homogeneous dispersion symbol of degree $m>1$, we prove a strict mass gap for every negative-energy traveling profile, uniformly over two unbounded classes of added dispersion and nonlinear energy. The dispersion addition satisfies $b\ge0$ and $\xi\cdot\nabla b\le mb$; no lower dilation bound is required. The nonlinear addition satisfies $G\ge0$ and $0\le\rho G'\le(1+m/N)G$. A shifted-frequency identity preserves favorable dilation defects even when their sizes are not uniformly bounded. Conversely, an arbitrarily small nonnegative relative-form perturbation of the same order can admit negative-energy stationary waves at every mass in an interval immediately above the original critical mass. Thus smallness alone does not preserve the profile gap. For the homogeneous flow we also prove a quantitative residual inequality with positive inertial response and a necessary lower bound on shape evolution for relative-periodic states. An explicit two-packet construction identifies an obstruction to a proposed pointwise interaction-virial sign. The arguments do not establish an unrestricted dynamical mass gap or classify all nonradial solutions.
\end{abstract}
\noindent\textbf{Status.} This manuscript contains analytic proofs with internal independent checks. External peer review and exhaustive publication-priority verification have not been completed. The mass-gap constants are qualitative. No numerical profile or computer-assisted certificate is used in these arguments.

\section*{The conclusions and their scope}
The reference equation is
\[
 i u_t=a(D)u-|u|^{2m/N}u,
\]
where $a$ is real, even, smooth away from zero, homogeneous of degree $m>1$, and comparable to $|\xi|^m$. Pure mass-critical BNLS is $a(\xi)=|\xi|^4$. Mass means $M(u)=\|u\|_2^2$, and a traveling wave has convention $u(t,x)=e^{i\omega t}U(x-vt)$.

The principal result is a uniform exclusion theorem for \emph{all traveling profiles} in the stated mass window, including nonradial profiles and arbitrary velocities. Its uniformity is stronger than perturbation by small coefficients: a single gap covers the entire admissible addition classes. The proof uses the sign of dilation defects, together with variational phase coercivity near the critical optimizer set.

The converse keeps the order of the operator unchanged. For any $\epsilon>0$ and $q>0$, put
\[
 b_{\epsilon,q}(\xi)=\epsilon a(\xi)\frac{|\xi|^q}{1+|\xi|^q},
 \qquad 0\le b_{\epsilon,q}(\xi)\le\epsilon a(\xi).
\]
The equation with dispersion $a+b_{\epsilon,q}$ has real negative-energy stationary profiles for every
\[
 M_{c,a}<M<(1+\epsilon)^{N/m}M_{c,a}.
\]
This construction lies in the same energy space and can be arbitrarily small in relative quadratic-form norm. The full modified kinetic form has the same sharp GN mass $M_{c,a}$, although its sharp inequality is no longer attained. The example violates the upper dilation inequality. It proves that this inequality cannot simply be discarded in favor of positivity and smallness; it does not claim a necessary-and-sufficient classification of every individual symbol.

For general time-dependent solutions, the residual theorem measures the failure to be a traveling profile. Its relative-periodic consequence forces nonzero shape motion; it does not rule out every possible oscillatory state. The explicit virial construction in Appendix~\ref{sec:r9-virial-obstruction} concerns initial derivatives and uniform estimates, not the later dynamics of its initial functions.

\section{A uniform mass gap for positive subhomogeneous additions}
\label{sec:cone}

Fix \(N\ge1\), \(m>1\), and a real-even symbol
\(a\in C^{N+2}(\mathbb R^N\setminus\{0\})\) such that
\begin{equation}\label{cone:principal}
a(t\xi)=t^m a(\xi),\qquad
c|\xi|^m\le a(\xi)\le C|\xi|^m
\end{equation}
for \(t>0\) and positive constants \(c,C\). Put
\[
s=\frac m2,\qquad r=\frac mN,\qquad p=2+2r,
\qquad A_a(f)=\int a(\xi)|\widehat f(\xi)|^2\,d\xi.
\]
The Fourier transform is unitary. Define \(M(f)=\|f\|_2^2\) and let
\(M_{c,a}\) be the sharp critical mass determined by
\begin{equation}\label{cone:gn}
\|f\|_p^p\le\frac p2
\left(\frac{M(f)}{M_{c,a}}\right)^r A_a(f),
\qquad f\in H^{m/2}(\mathbb R^N;\mathbb C).
\end{equation}
The reference mass in this section belongs to the principal
homogeneous symbol \(a\).

\subsection{The two admissible classes}

An added dispersion symbol \(b\) is admissible if it is real and even,
belongs to \(C^{N+2}(\mathbb R^N\setminus\{0\})\), and satisfies
\begin{align}
b(\xi)&\ge0,\qquad \xi\cdot\nabla b(\xi)\le m b(\xi),
&&\xi\ne0,\label{cone:dispersion-sign}\\
|\partial^\beta b(\xi)|
&\le C_{b,\beta}|\xi|^{-|\beta|}(1+|\xi|^m),
&&|\beta|\le N+2,\quad \xi\ne0.\label{cone:dispersion-regularity}
\end{align}
Each constant \(C_{b,\beta}\) must be finite for its individual symbol;
there is no uniform bound on these constants across the class.
There is no lower bound on \(\xi\cdot\nabla b\). Values at the single
frequency zero do not affect the multiplier.

An added nonlinear energy density \(G\) is admissible if
\begin{equation}\label{cone:nonlinear-sign}
G\in C^1([0,\infty)),\quad G(0)=0,\quad G\ge0,\qquad
0\le\rho G'(\rho)\le(1+r)G(\rho)\quad(\rho>0),
\end{equation}
and the map
\begin{equation}\label{cone:nonlinear-regularity}
z\longmapsto g(|z|^2)z,\qquad g=2G',
\end{equation}
is locally \(C^1\) on \(\mathbb C\simeq\mathbb R^2\).
This last requirement is separate from the inequalities in
\eqref{cone:nonlinear-sign}; it supplies the Sobolev chain rule.
These assumptions imply, for a finite constant depending on \(G\),
\begin{equation}\label{cone:nonlinear-growth}
0\le G(\rho)\le C_G(\rho+\rho^{1+r}),\qquad
0\le g(\rho)\le C_G(1+\rho^r).
\end{equation}
For \(\rho\le1\) this follows from the continuity of \(G'\) and
\(G(0)=0\). For \(\rho\ge1\), integrate
\(\rho G'(\rho)\le(1+r)G(\rho)\), or apply the differential inequality
directly if \(G(1)=0\), and then use it to bound \(G'\).

Write
\[
B_b(f)=\int b(\xi)|\widehat f(\xi)|^2\,d\xi,\qquad
F_G(f)=\int G(|f|^2)\,dx,
\]
and define
\begin{equation}\label{cone:energy}
E_{a,b,G}(f)=\frac12\bigl(A_a(f)+B_b(f)\bigr)
-\frac1p\|f\|_p^p+F_G(f).
\end{equation}
All these expressions are finite on \(H^{m/2}\) by
\eqref{cone:dispersion-regularity},
\eqref{cone:nonlinear-growth}, and Sobolev interpolation.

\begin{theorem}[Uniform mass gap over the full addition classes]
\label{cone:gap}
There exists \(\varepsilon_{N,m,a}>0\) such that, for every admissible
\(b,G\), every \(\omega\in\mathbb R\), and every \(v\in\mathbb R^N\),
each nonzero weak \(H^{m/2}\) solution of
\begin{equation}\label{cone:profile}
(a(D)+b(D))U+\omega U+iv\cdot\nabla U
-|U|^{2r}U+g(|U|^2)U=0
\end{equation}
with \(E_{a,b,G}(U)<0\) satisfies
\[
M(U)\ge M_{c,a}+\varepsilon_{N,m,a}.
\]
The gap is independent of the coefficients, derivative bounds,
and detailed shapes of the admissible additions. It is also
independent of frequency and velocity.
\end{theorem}

The theorem concerns traveling profiles. It does not assert their
existence for every admissible model, nor does it assert a
decomposition or classification of arbitrary global solutions.
The proof does not require radiality, finite spatial moments,
uniqueness of the principal optimizer, or real-sector
nondegeneracy of that optimizer.

\subsection{Regularity and the cutoff dilation identity}

\begin{lemma}[Individual profile regularity]\label{cone:regularity}
Every weak \(H^{m/2}\) profile in \eqref{cone:profile} belongs to
\(L^\infty\cap H^{m+1}\).
\end{lemma}
\begin{proof}
For the fixed symbol \(b\) and velocity \(v\), choose \(\Lambda>0\)
large enough that
\[
a(\xi)+b(\xi)-v\cdot\xi+\Lambda\asymp1+|\xi|^m.
\]
Here \(b\ge0\), ellipticity, and \(m>1\) give the lower bound,
and \eqref{cone:dispersion-regularity} gives the upper bound.
The resolvent
\[
R_\Lambda=(a(D)+b(D)+iv\cdot\nabla+\Lambda)^{-1}
\]
maps \(L^q\) into the Bessel potential space \(H^{m,q}\) for every
\(1<q<\infty\). Indeed the multiplier
\[
\frac{\langle\xi\rangle^m}
{a(\xi)+b(\xi)-v\cdot\xi+\Lambda}
\]
satisfies the Mihlin bounds by
\eqref{cone:dispersion-regularity} and homogeneity of \(a\).
Constants may depend on the individual model and velocity.

The equation takes the form
\[
U=R_\Lambda\bigl(|U|^{2r}U-g(|U|^2)U\bigr)
+(\Lambda-\omega)R_\Lambda U.
\]
The nonlinear expression has magnitude at most
\(C(|U|+|U|^{1+2r})\) by \eqref{cone:nonlinear-growth}.
For completeness, when \(N>m\), start with
\(q_0=2N/(N-m)\) and reciprocal exponent
\(\beta_0=1/2-m/(2N)\). On \(\{|U|>1\}\) the nonlinear
expression is bounded by \(C|U|^{1+2r}\), and on
\(\{|U|\le1\}\) by \(C|U|\). Apply the resolvent to these
two terms separately. Sobolev embedding and interpolation give
the improvement
\[
\beta_{j+1}=(1+2r)\beta_j-r
\]
until an exponent exceeds the threshold for boundedness. The
linear contribution has the better reciprocal output exponent
\(\beta_j-r\) and remains in \(L^2\); interpolation gives a common
output exponent for the two terms. The ideal iterates are
\(\beta_j=1/2-(r/2)(1+2r)^j\), so finitely many iterations
suffice. All nonlinear input exponents exceed one since
\(q_0>1+2r\), and intermediate exponents can be weakened slightly
to avoid endpoints. If \(N=m\), the initial energy space embeds
in every finite \(L^q\), and one starts with sufficiently large
\(q\); if \(N<m\), it already embeds into \(L^\infty\).
Thus \(U\in L^\infty\). The right-hand nonlinearities are now in
\(L^2\), giving \(U\in H^m\subset H^1\). The local \(C^1\)
assumption \eqref{cone:nonlinear-regularity} and the corresponding
property of \(z\mapsto|z|^{2r}z\) imply that both nonlinearities
belong to \(H^1\). A final resolvent application gives
\(U\in H^{m+1}\).
\end{proof}

Choose \(\chi\in C_c^\infty(\mathbb R^N;\mathbb R)\) equal to one
near zero, and put
\[
f_R(x)=x\chi(x/R)=R f(x/R),\qquad
L_R=f_R\cdot\nabla+\frac12\operatorname{div}f_R,\qquad R\ge1.
\]
The operator \(L_R\) is skew-symmetric on \(L^2\).

\begin{lemma}[Nonhomogeneous cutoff commutator]\label{cone:commutator}
The commutators \([b(D),L_R]\) are uniformly bounded in \(R\)
as forms \(H^{m/2}\to H^{-m/2}\). For every \(u\in H^{m/2}\),
\begin{equation}\label{cone:commutator-limit}
\langle u,[b(D),L_R]u\rangle
\longrightarrow
\int \xi\cdot\nabla b(\xi)|\widehat u(\xi)|^2\,d\xi.
\end{equation}
The bound may depend on the individual symbol \(b\).
\end{lemma}
\begin{proof}
Apart from the fixed Fourier normalization, the kernel is
\[
\frac i2\bigl(b(\xi)-b(\eta)\bigr)
(\xi+\eta)\cdot\widehat f_R(\xi-\eta).
\]
The zeroth and first derivative bounds in
\eqref{cone:dispersion-regularity} imply
\begin{equation}\label{cone:kernel-estimate}
|b(\xi)-b(\eta)|\,|\xi+\eta|
\le C_b|\xi-\eta|\bigl(1+(|\xi|+|\eta|)^m\bigr).
\end{equation}
To see this, write \(\rho=|\xi|+|\eta|\). If
\(|\xi-\eta|\ge\rho/4\), use the zeroth-order bound.
Otherwise the segment between the two frequencies has magnitude
comparable to \(\rho\), and the first derivative bound and
mean-value theorem apply.

For \(k=\xi-\eta\), divide \eqref{cone:kernel-estimate} by
\(\langle\xi\rangle^{m/2}\langle\eta\rangle^{m/2}\).
The result is bounded by \(C_b|k|\langle k\rangle^{m/2}\);
this follows from the inequalities for
\(\langle\xi\rangle/\langle\eta\rangle\) and its reciprocal
by a constant times \(\langle k\rangle\).
Since \(\widehat f_R(k)=R^{N+1}\widehat f(Rk)\),
\[
\int |k|\langle k\rangle^{m/2}|\widehat f_R(k)|\,dk
=\int |\ell|\langle\ell/R\rangle^{m/2}
|\widehat f(\ell)|\,d\ell\le C_f.
\]
Young's inequality proves the uniform form bound.

For Schwartz \(u\), \(L_Ru\to Lu\) in \(L^2\), where
\(L=x\cdot\nabla+N/2\). Thus the left side of
\eqref{cone:commutator-limit} converges to
\(2\operatorname{Re}\langle b(D)u,Lu\rangle\).
Fourier integration by parts gives the asserted integral. One may
first exclude a ball of radius \(\delta\) around zero: \(b\)
is bounded there by \eqref{cone:dispersion-regularity}, so the
boundary term is \(O(\delta^N)\); the term at infinity vanishes
by Schwartz decay. Finally
\[
|\xi\cdot\nabla b(\xi)|\le C_b(1+|\xi|^m)
\]
by the same derivative assumption. Density and the uniform form
bound extend the limit to \(H^{m/2}\). This argument does not use
a lower radial-slope bound.
\end{proof}

The identical argument applies to \(a\), giving limit \(mA_a(u)\)
by homogeneity. Applied to \(iv\cdot\nabla\), it gives limit
\(-v\cdot P(u)\), where
\[
P(u)=\int \xi|\widehat u(\xi)|^2\,d\xi.
\]
This momentum is finite because \(m>1\). The degree-one
commutator is bounded on \(H^{1/2}\), and therefore on the
energy space.

\subsection{Exact defects and the shifted frequency}

For a profile \(U\), abbreviate
\[
M=M(U),\quad A=A_a(U),\quad B=B_b(U),\quad
F=F_G(U),\quad E=E_{a,b,G}(U),\quad \rho=|U|^2.
\]
Define the two dilation defects
\begin{align}
D_b&=\int\bigl(mb(\xi)-\xi\cdot\nabla b(\xi)\bigr)
|\widehat U(\xi)|^2\,d\xi\ge0,\label{cone:defect-b}\\
K_G&=N\int\bigl(\rho G'(\rho)-G(\rho)\bigr)\,dx,\qquad
D_G=mF-K_G\ge0.\label{cone:defect-G}
\end{align}
Each defect is finite. In fact
\(-NF\le K_G\le mF\), so \(0\le D_G\le(m+N)F\).
There is no required estimate of \(D_b\) by \(B\), uniformly
over the dispersion class.

\begin{lemma}[Profile identities]\label{cone:identities}
Every profile obeys
\begin{align}
v\cdot P(U)&=2mE-D_b-2D_G,\label{cone:pohozaev}\\
\omega M&=r(A+B)+(2m-p)E-D_b-\vartheta_ND_G,
\qquad \vartheta_N=2(1-1/N).\label{cone:frequency}
\end{align}
Writing \(U=X+iY\) with real \(X,Y\), one has
\begin{equation}\label{cone:phase-energy}
mE=\left\langle Y,
\bigl[a(D)+b(D)+\omega-|U|^{2r}+g(|U|^2)\bigr]Y\right\rangle
+\frac{D_b}{2}+D_G.
\end{equation}
\end{lemma}
\begin{proof}
Lemma~\ref{cone:regularity} makes \(L_RU\) an admissible test.
The kinetic and drift terms converge by
Lemma~\ref{cone:commutator}. The frequency pairing vanishes.
The critical focusing term converges to
\((m/p)\|U\|_p^p\). For the added nonlinearity, the chain rule gives
\(\nabla G(|U|^2)=g(|U|^2)\operatorname{Re}(\overline U\nabla U)\),
and compact-support integration by parts yields
\[
\operatorname{Re}\int g(\rho)U\,\overline{L_RU}\,dx
=\int\operatorname{div}f_R\bigl(\rho G'(\rho)-G(\rho)\bigr)\,dx
\longrightarrow K_G.
\]
Its integrand is dominated by a constant times \(G(\rho)\).
These limits give
\[
\frac12\left(mA+\int\xi\cdot\nabla b\,|\widehat U|^2\right)
-\frac12v\cdot P(U)-\frac mp\|U\|_p^p+K_G=0,
\]
which is \eqref{cone:pohozaev}.

Pairing the profile with \(U\) gives
\[
A+B+\omega M-v\cdot P(U)-\|U\|_p^p+J=0,
\qquad J=\int g(\rho)\rho\,dx.
\]
Since
\[
J=2F+\frac2N K_G,\qquad
pF-J=\frac2N D_G,\qquad
\|U\|_p^p=\frac p2(A+B)+pF-pE,
\]
substitution proves \eqref{cone:frequency}.

Real-even symbols preserve real functions, so the imaginary
profile equation reads
\[
\bigl[a(D)+b(D)+\omega-|U|^{2r}+g(|U|^2)\bigr]Y
=-v\cdot\nabla X.
\]
Also \(P(U)=-2\int Y\nabla X\). Pairing this equation with \(Y\)
and using \eqref{cone:pohozaev} proves
\eqref{cone:phase-energy}.
\end{proof}

The useful frequency is
\begin{equation}\label{cone:shifted-frequency}
\widetilde\omega=\omega+\frac{D_b+\vartheta_ND_G}{M},
\qquad
\widetilde\omega M=r(A+B)+(2m-p)E.
\end{equation}
Substituting this frequency into \eqref{cone:phase-energy} gives
the exact identity
\begin{equation}\label{cone:shifted-identity}
\begin{aligned}
mE={}&
\left\langle Y,\bigl[a(D)+\widetilde\omega-|U|^{2r}\bigr]Y\right\rangle
+B_b(Y)+\int g(|U|^2)Y^2\,dx\\
&+D_b\left(\frac12-\frac{M(Y)}M\right)
+D_G\left(1-\vartheta_N\frac{M(Y)}M\right).
\end{aligned}
\end{equation}
This identity retains the compensating contributions of the
defects. In particular, an arbitrarily large negative shift of
the actual frequency need not be estimated separately.

\subsection{Proof of the uniform gap}

\begin{proof}[Proof of Theorem~\ref{cone:gap}]
The sharp inequality and positivity of \(B,F\) imply that
negative energy requires \(M>M_{c,a}\). If the asserted gap
failed, choose admissible additions \(b_n,G_n\) and profiles
\(U_n\) such that \(E_n<0\) and \(M(U_n)\to M_{c,a}\).

Critical dilation acts on the entire model by
\begin{align*}
U_\lambda(x)&=\lambda^{N/2}U(\lambda x),&
b_\lambda(\xi)&=\lambda^m b(\xi/\lambda),\\
G_\lambda(\rho)&=\lambda^{m+N}G(\lambda^{-N}\rho),&
g_\lambda(\rho)&=\lambda^m g(\lambda^{-N}\rho),\\
\omega_\lambda&=\lambda^m\omega,&
v_\lambda&=\lambda^{m-1}v.
\end{align*}
It preserves both admissible classes and the profile equation,
preserves mass, and multiplies the energy by \(\lambda^m\).
Thus normalize \(A_a(U_n)=1\). The individual regularity constants
may change under scaling, but no uniform bound on them is used.

Set
\[
\delta_n=\left(\frac{M(U_n)}{M_{c,a}}\right)^r-1\longrightarrow0.
\]
The sharp inequality and negative energy give
\begin{equation}\label{cone:near-threshold}
0\le B_n+2F_n<\delta_n,\qquad
-\frac{\delta_n}{2}\le E_n<0,\qquad
\|U_n\|_p^p\longrightarrow\frac p2.
\end{equation}
Lemma~\ref{foundation:compact} gives translations and constant
phases for which \(U_n\to Q\) strongly in \(H^{m/2}\), where
\(Q\) is a real optimizer normalized by
\[
M(Q)=M_{c,a},\qquad A_a(Q)=1,\qquad
\omega_Q=r/M_{c,a}.
\]
Equation \eqref{cone:shifted-frequency} gives
\(\widetilde\omega_n\to\omega_Q\).
After a further constant phase tending to zero, arrange
\(\langle Y_n,Q\rangle=0\), where \(Y_n=\operatorname{Im}U_n\).

By Lemma~\ref{foundation:phase}, the principal phase form
\(L_-=a(D)+\omega_Q-|Q|^{2r}\) is coercive on real functions
\(L^2\)-orthogonal to \(Q\). Strong \(L^p\) convergence gives
\[
|U_n|^{2r}\longrightarrow|Q|^{2r}
\quad\hbox{in }L^{p/(2r)}.
\]
H\"older's inequality with two \(L^p\) factors and Sobolev
embedding transfer this coercivity to the shifted forms:
\[
\left\langle Y_n,
\bigl[a(D)+\widetilde\omega_n-|U_n|^{2r}\bigr]Y_n\right\rangle
\ge\frac{c_Q}{2}\|Y_n\|_{H^{m/2}}^2
\]
for all sufficiently large \(n\).

Moreover \(M(Y_n)/M(U_n)\to0\). Once this ratio is at most \(1/4\),
and since \(0\le\vartheta_N<2\), identity
\eqref{cone:shifted-identity} gives
\[
\begin{aligned}
mE_n\ge{}&
\frac{c_Q}{2}\|Y_n\|_{H^{m/2}}^2+B_{b_n}(Y_n)
+\int g_n(|U_n|^2)Y_n^2\,dx\\
&+\frac14D_{b,n}+\frac12D_{G,n}\ge0.
\end{aligned}
\]
All displayed terms are finite and nonnegative; for the
multiplication term use
\(\int g_n(|U_n|^2)Y_n^2\le J_n\le2(1+r)F_n\).
This contradicts \(E_n<0\). No boundedness of the actual
frequencies, and no convergence or boundedness of \(D_{b,n}\),
was needed.
\end{proof}

\subsection{Uniformity under changes of the principal symbol}

\begin{theorem}[Local uniform robustness]\label{cone:robustness}
Fix a principal symbol \(a_0\) satisfying
\eqref{cone:principal}. There are \(\eta>0\) and
\(\varepsilon>0\) such that the conclusion
\[
M(U)\ge M_{c,a}+\varepsilon
\]
holds for every individually admissible principal symbol \(a\)
with
\[
\sup_{\xi\in S^{N-1}}\left|\frac{a(\xi)}{a_0(\xi)}-1\right|<\eta,
\]
simultaneously for every admissible pair \(b,G\).
No uniform derivative bounds on any symbol are required.
\end{theorem}
\begin{proof}
If false, choose \(a_n\to a_0\) uniformly relative to \(a_0\)
on the sphere, admissible additions, and negative-energy profiles
whose masses differ from \(M_{c,a_n}\) by a quantity tending to
zero. The principal form comparison
\[
(1-o(1))A_{a_0}\le A_{a_n}\le(1+o(1))A_{a_0}
\]
implies convergence of sharp constants and hence
\(M_{c,a_n}\to M_{c,a_0}\). Normalize \(A_{a_n}(U_n)=1\).
The bounds \eqref{cone:near-threshold} hold with \(M_{c,a_n}\);
in addition \(A_{a_0}(U_n)\to1\). Variational compactness for
\(a_0\) gives convergence to its real normalized optimizer \(Q\).
The shifted frequency again converges to \(r/M_{c,a_0}\).
The principal forms converge in norm from \(H^{m/2}\) to
\(H^{-m/2}\), so their phase coercivity transfers exactly as in
the preceding proof. The added forms and dilation defects are
retained as nonnegative terms in
\eqref{cone:shifted-identity}. The same contradiction follows.
\end{proof}

\begin{proposition}[Compact principal families]\label{cone:compact-family}
Let the restrictions to \(S^{N-1}\) of a family of principal
symbols form a compact subset of \(C^0(S^{N-1})\), with every
member satisfying \eqref{cone:principal}. A single positive gap
works over that family and over both full addition classes.
\end{proposition}
\begin{proof}
Apply Theorem~\ref{cone:robustness} at every member, extract
a finite subcover, and take the minimum of the corresponding
positive gap constants.
\end{proof}

\subsection{Examples and the structural mechanism}

For \(a(\xi)=|\xi|^4\), the theorem includes
\(b(\xi)=\mu|\xi|^2\) for every \(\mu\ge0\), with one common
gap independent of \(\mu\). More generally it includes arbitrary
finite nonnegative sums of real-even nonnegative homogeneous
symbols of orders \(0<k\le m\), and nonnegative constant terms.
The added symbols need not be elliptic.

Nonhomogeneous examples include
\[
b(\xi)=\mu\frac{|\xi|^k}{1+|\xi|^k}\quad(0<k\le m),\qquad
b(\xi)=\mu e^{-|\xi|^2},\qquad
b(\xi)=\mu(1+|\xi|^2)^{-q},
\]
simultaneously for arbitrary \(\mu,q\ge0\). The Gaussian has
logarithmic radial slope \(-2|\xi|^2\), which is unbounded below;
this illustrates why the shifted-frequency argument is stronger
than imposing a two-sided radial-slope estimate.
Nonnegative smooth radially decreasing compactly supported
frequency multipliers are also admissible.

The nonlinear class includes all finite sums
\[
G(\rho)=\sum_j\frac{\kappa_j}{q_j+2}\rho^{1+q_j/2},
\qquad \kappa_j\ge0,\qquad 0\le q_j\le2m/N,
\]
which add \(\sum_j\kappa_j|U|^{q_j}U\) to the profile equation.
It also includes, for arbitrary \(\kappa\ge0\),
\[
G(\rho)=\kappa\frac{\rho^{1+r}}{(1+\rho)^r},\qquad
G(\rho)=\kappa\log(1+\rho),\qquad
G(\rho)=\kappa\frac{\rho}{1+\rho}.
\]
The latter two are concave energy densities. Thus neither
convexity of \(G\) nor monotonicity of \(g=2G'\) is required.

The argument uses two complementary consequences of being near
the sharp principal mass. Negative energy forces the actual added
energies \(B+2F\) to be small after kinetic normalization, even if
their coefficients are arbitrarily large. Variational compactness
then makes the nonconstant phase component small. In the shifted
identity, that small phase mass allows the positive dilation
defects to dominate their own negative contributions to the
frequency. This is a sign-and-scaling mechanism rather than a
small-operator perturbation argument.

\section{Arbitrarily small nonhomogeneous perturbations can remove the gap}
\label{sec:r9-bounded-boundary}

The upper dilation condition has a substantive role even among symbols
of the same order and arbitrarily close in quadratic-form norm. The
following construction complements the exclusion theorem. Its existence
proof uses standard constrained minimization; no separate priority claim
is attached to that variational method.

Fix a principal symbol $a$ as in the cone theorem, of degree $m>1$,
and write
\[
 r=m/N,\qquad p=2+2r,\qquad M_c=M_{c,a},\qquad
 A(f)=\int a(\xi)|\widehat f(\xi)|^2\,d\xi.
\]
For fixed numbers $\varepsilon,q>0$, set
\begin{equation}\label{eq:r9-bounded-symbol}
 b_{\varepsilon,q}(\xi)
 =\varepsilon a(\xi)\frac{|\xi|^q}{1+|\xi|^q},\qquad
 B(f)=\int b_{\varepsilon,q}(\xi)|\widehat f(\xi)|^2\,d\xi,
 \qquad E(f)=\frac12[A(f)+B(f)]-\frac1p\|f\|_p^p.
\end{equation}
The symbol is real, even, individually $C^{N+2}$ away from zero, and
satisfies all the order-$m$ derivative bounds of the cone theorem. In
particular,
\begin{equation}\label{eq:r9-bounded-formsmall}
 0\le b_{\varepsilon,q}\le\varepsilon a,
 \qquad 0\le B(f)\le\varepsilon A(f).
\end{equation}
But its dilation defect has the opposite sign:
\begin{equation}\label{eq:r9-bounded-defect}
 mb_{\varepsilon,q}-\xi\cdot\nabla b_{\varepsilon,q}
 =-\varepsilon q\,a(\xi)
       \frac{|\xi|^q}{(1+|\xi|^q)^2}<0\quad(\xi\ne0).
\end{equation}

The number $M_c$ is also the actual sharp GN threshold for the full
kinetic form $A+B$, not merely a comparison with the principal part.
Indeed $A+B\ge A$ gives the principal sharp constant as an upper
bound for the full one. If $Q$ is a principal optimizer and
$Q_\lambda(x)=\lambda^{N/2}Q(\lambda x)$, dominated convergence gives
\[
 \frac{B(Q_\lambda)}{\lambda^m}
 =\varepsilon\int a(\xi)
 \frac{\lambda^q|\xi|^q}{1+\lambda^q|\xi|^q}
 |\widehat Q(\xi)|^2\,d\xi\longrightarrow0.
\]
The sharp constants therefore coincide. The full inequality has no
nonzero optimizer: $B(f)>0$ makes its comparison with the principal
sharp inequality strict for every nonzero $f$.

\begin{theorem}[A boundary within the same operator order]
\label{thm:r9-bounded-boundary}
Define
\[
 I(M)=\inf\{E(f): f\in H^{m/2}(\mathbb R^N;\mathbb C),\ \|f\|_2^2=M\}.
\]
For $0<M\le M_c$, $I(M)=0$ and the infimum is not attained.
For every
\begin{equation}\label{eq:r9-bounded-window}
 M_c<M<(1+\varepsilon)^{N/m}M_c,
\end{equation}
$-\infty<I(M)<0$, and each minimizing sequence has, after translations
and passage to a subsequence, a strong $H^{m/2}$ limit attaining the
infimum. A minimizer $Q_M$ can be chosen real and belongs to
$L^\infty\cap H^{m+1}$. For some $\omega_M>0$ it solves
\begin{equation}\label{eq:r9-bounded-profile}
 [a(D)+b_{\varepsilon,q}(D)]Q_M+\omega_M Q_M
       -|Q_M|^{2r}Q_M=0.
\end{equation}
Its energy is exactly
\begin{equation}\label{eq:r9-bounded-energy}
 E(Q_M)=-\frac{\varepsilon q}{2m}
 \int a(\xi)\frac{|\xi|^q}{(1+|\xi|^q)^2}
             |\widehat Q_M(\xi)|^2\,d\xi<0.
\end{equation}
Thus the modified evolution has a real-profile, zero-momentum standing
wave $u_M(t)=e^{i\omega_Mt}Q_M$ for every mass in
\eqref{eq:r9-bounded-window}. For each fixed $\varepsilon>0$, these
negative-energy standing-wave masses approach $M_c$ from above.
\end{theorem}

\begin{proof}
The sharp principal inequality is
\begin{equation}\label{eq:r9-bounded-GN}
 \frac1p\|f\|_p^p\le
 \frac12(M(f)/M_c)^r A(f).
\end{equation}
Consequently, if $0<M\le M_c$ then $E(f)\ge B(f)/2>0$ for every
nonzero admissible $f$. The strict inequality holds because
$b_{\varepsilon,q}>0$ away from the single frequency zero. A smooth
mass-$M$ trial function dilated to low frequency has energy tending to
zero, so $I(M)=0$ without attainment.

Now fix $M$ in \eqref{eq:r9-bounded-window}, and put
$d=(M/M_c)^r-1\in(0,\varepsilon)$. For every
$\eta\in(0,\varepsilon)$ there is a finite $C_\eta$ such that
\[
 b_{\varepsilon,q}(\xi)\ge(\varepsilon-\eta)a(\xi)-C_\eta.
\]
Indeed the coefficient $\varepsilon|\xi|^q/(1+|\xi|^q)$ tends to
$\varepsilon$ at infinity, and $a$ is bounded on every finite ball.
Choose $\eta<\varepsilon-d$. Equation~\eqref{eq:r9-bounded-GN} gives
\begin{equation}\label{eq:r9-bounded-coercivity}
 E(f)\ge\frac{\varepsilon-\eta-d}{2}A(f)-\frac{C_\eta M}{2}
 \quad\text{when }M(f)=M.
\end{equation}
Thus $I(M)>-\infty$, and minimizing sequences are bounded in
$H^{m/2}$.

To see strict negativity, multiply a real principal GN optimizer of
mass $M_c$ by $\sqrt{M/M_c}$ and approximate it in $H^{m/2}$ by real
$C_c^\infty$ functions, adjusting the mass. This yields a real smooth
$f$ of mass exactly $M$ with
$E_a(f)=A(f)/2-\|f\|_p^p/p<0$.
For $f_\lambda(x)=\lambda^{N/2}f(\lambda x)$,
\[
 B(f_\lambda)\le\varepsilon\lambda^{m+q}
       \int a(\xi)|\xi|^q|\widehat f(\xi)|^2\,d\xi,
 \qquad E(f_\lambda)=\lambda^m E_a(f)+\tfrac12B(f_\lambda)<0
\]
for sufficiently small positive $\lambda$. This proves $I(M)<0$.
The same statements hold at every smaller positive mass in the window
where needed below.

We give the compactness argument to handle the mass-critical power
without presuming that the added symbol has higher order.
For $\theta>1$ such that the masses under consideration remain below
the upper endpoint of \eqref{eq:r9-bounded-window}, amplitude scaling
gives
\[
 I(\theta L)\le\theta I(L),\qquad
 I(\theta L)\le\theta^{p/2}I(L)<\theta I(L)
       \quad\text{if }I(L)<0.
\]
The second inequality follows from
\[
 E(\sqrt\theta f)=\theta^{p/2}E(f)
       +\tfrac12(\theta-\theta^{p/2})[A(f)+B(f)].
\]
For every split $M=M_1+M_2$ with $M_j>0$, this implies
\begin{equation}\label{eq:r9-bounded-binding}
 I(M)<I(M_1)+I(M_2).
\end{equation}
If both pieces have nonnegative infimum, their infima are zero while
$I(M)<0$. Otherwise $I(M)/M$ is strictly below the energy-per-mass
ratio of a negative piece and no greater than that of the other piece;
weighting by $M_1,M_2$ proves \eqref{eq:r9-bounded-binding}.

A bounded minimizing sequence $f_n$ cannot vanish locally in mass.
Bounded-overlap localization of the GN inequality gives
\[
 \|f_n\|_p^p
 \le C\left(\sup_y\int_{B(y,2)}|f_n|^2\right)^r
           \|f_n\|_{H^{m/2}}^2.
\]
For example, apply the GN inequality to smooth unit-cube cutoffs and
sum their squared $H^{m/2}$ norms, which are bounded by a constant
times the global squared norm. Local vanishing would imply
$\|f_n\|_p\to0$, contradicting $I(M)<0$.
After translations and extraction, there is therefore a nonzero weak
limit $Q$ in $H^{m/2}$, with almost-everywhere convergence. Put
$w_n=f_n-Q$. The two quadratic forms split by weak Hilbert convergence
(the $b$ form is bounded by $\varepsilon A$), and the nonlinear
integral splits by the Brezis--Lieb lemma. If $0<M(Q)<M$, normalizing
$w_n$ to mass exactly $M-M(Q)$ changes its energy by $o(1)$, because
its energy terms are bounded. The splitting would then contradict
\eqref{eq:r9-bounded-binding}. Hence $M(Q)=M$ and $M(w_n)\to0$.
The principal GN inequality gives
\[
 \|w_n\|_p^p\le C M(w_n)^r A(w_n)\longrightarrow0.
\]
Energy splitting and $E(Q)\ge I(M)$ now force
$A(w_n)+B(w_n)\to0$. This proves strong $H^{m/2}$ convergence and
attainment.

A real minimizer is available without any modulus inequality for
higher-order kinetic forms. Write $f=X+iY$ and, when both real
components have positive mass, set
$\vartheta=M(X)/M(f)$, $g=X/\sqrt\vartheta$ and
$h=Y/\sqrt{1-\vartheta}$. Both $g,h$ have mass $M(f)$. Real-evenness
splits the kinetic form, whereas convexity gives
\[
 |f|^p=[\vartheta g^2+(1-\vartheta)h^2]^{p/2}
       \le\vartheta|g|^p+(1-\vartheta)|h|^p.
\]
Thus $E(f)\ge\vartheta E(g)+(1-\vartheta)E(h)$, and the complex
and real constrained infima coincide. A real minimizing sequence
and the proved compactness yield a real minimizer.

The Lagrange multiplier rule gives \eqref{eq:r9-bounded-profile}.
Testing it against $Q_M$ shows
\[
 \omega_M M=\|Q_M\|_p^p-A(Q_M)-B(Q_M)
       =\left(1-\frac2p\right)\|Q_M\|_p^p-2E(Q_M)>0.
\]
The symbol $a+b_{\varepsilon,q}+\omega_M$ is comparable with
$1+|\xi|^m$ and its resolvent satisfies the individual Mihlin bounds.
The finite Sobolev bootstrap for the principal critical power therefore
gives $L^\infty$, then $H^m\subset H^1$, and finally $H^{m+1}$ by
the local $C^1$ Sobolev chain rule. This is the same individual
regularity argument as for the cone theorem and uses no sign of the
dilation defect.

Finally the mass-preserving scaling curve through the minimizer has
its scalar energy minimum at $\lambda=1$. Differentiation of its
explicit Fourier energy formula is justified by
$|\xi\cdot\nabla b_{\varepsilon,q}|\le(m+q)b_{\varepsilon,q}$
and domination by $a$. It yields
\[
 \frac m2 A(Q_M)
 +\frac12\int\xi\cdot\nabla b_{\varepsilon,q}(\xi)
                         |\widehat Q_M(\xi)|^2\,d\xi
 -\frac mp\|Q_M\|_p^p=0.
\]
Subtracting this identity divided by $m$ from the energy formula
proves \eqref{eq:r9-bounded-energy}. No spatial moment or
differentiability of the scaling curve in the energy norm is needed.
A real $Q_M$ has zero momentum, and its standing-wave formula supplies
a global solution directly.
\end{proof}

For the original fourth-order principal symbol, one concrete example is
\[
 a(\xi)=|\xi|^4,\qquad
 b_{\varepsilon,2}(\xi)=\varepsilon\frac{|\xi|^6}{1+|\xi|^2}.
\]
Its perturbation is at most $\varepsilon$ times the original kinetic
form, yet its negative-energy zero-momentum standing waves have every
mass in $(M_c,(1+\varepsilon)^{N/4}M_c)$. This is a result about the
specified modified equation, not a counterexample for the pure symbol
$|\xi|^4$.

The construction does not say that every symbol violating the upper
dilation inequality loses its gap. It proves that removing that
condition admits counterexamples arbitrarily close to the principal
operator, even with the same energy space and the same order bounds.
In particular, the homogeneous-symbol robustness theorem cannot be
extended to all small nonhomogeneous quadratic-form perturbations.
The distinction is structural: smallness alone does not control the
sign of the dilation defect.

\section{Residual inertia and quantitative shape evolution}
\label{sec:r9-relative}

Let $N\ge1$, $m>1$, $s=m/2$, and let $a$ be a real, even,
positive elliptic homogeneous symbol of degree $m$, of class
$C^{N+2}$ away from the origin. Thus
$c|\xi|^m\le a(\xi)\le C|\xi|^m$.
Set $\alpha=2m/N$ and $p=2+2m/N$, and write
\[
 M(f)=\|f\|_2^2,\qquad
 A(f)=\int a(\xi)|\widehat f(\xi)|^2\,d\xi,\qquad
 E(f)=\tfrac12A(f)-\tfrac1p\|f\|_p^p,
 \qquad P(f)=\int\xi|\widehat f(\xi)|^2\,d\xi.
\]
We use the equivalent energy norm and its dual norm
\[
 \|f\|_X^2=M(f)+A(f),\qquad X=H^{m/2},\qquad
 \|h\|_{X^*}^2=\int\frac{|\widehat h(\xi)|^2}{1+a(\xi)}\,d\xi.
\]
Since $m>1$, differentiation maps $X$ continuously into $X^*$.
The map $f\mapsto |f|^\alpha f$ is locally Lipschitz from $X$ to
$X^*$, by H\"older's inequality and the embedding $X\hookrightarrow L^p$.
Pairings of real functions below are real dual pairings.

Let $M_c$ be the sharp critical mass, and let $\mathcal Q$ be the
set of real sharp optimizers normalized by $M(Q)=M_c$ and $A(Q)=1$.
Every $Q\in\mathcal Q$ satisfies
\[
 a(D)Q+\omega_*Q-|Q|^\alpha Q=0,
 \qquad \omega_*=\frac{m}{NM_c}.
\]
We use the optimizer compactness and phase coercivity established in
Lemmas~\ref{foundation:compact} and~\ref{foundation:phase}.
Their uniform consequences are recorded explicitly first.

\subsection{Uniform control near the optimizer set}

\begin{lemma}[Uniform optimizer bounds]\label{lem:r9rel-uniform}
There are $c_0,\kappa>0$, depending only on the fixed model, such that
for every $Q\in\mathcal Q$, the form
\[
 L_Q=a(D)+\omega_*-|Q|^\alpha
\]
obeys
\begin{align}
 \langle b,L_Qb\rangle&\ge c_0\|b\|_X^2
 &&\text{for real }b\in X,\ \langle b,Q\rangle=0,
 \label{eq:r9rel-uniformcoercivity}\\
 \|\zeta\cdot\nabla Q\|_{X^*}&\ge\kappa|\zeta|
 &&\text{for }\zeta\in\mathbb R^N,
 \label{eq:r9rel-derivativeinjective}\\
 \|\tau Q+i\zeta\cdot\nabla Q\|_{X^*}
 &\ge\kappa(|\tau|+|\zeta|)
 &&\text{for }(\tau,\zeta)\in\mathbb R\times\mathbb R^N.
 \label{eq:r9rel-parameterinjective}
\end{align}
\end{lemma}
\begin{proof}
The normalized optimizer set is compact in $X$ modulo translations.
Moreover, $Q\mapsto |Q|^\alpha$ is continuous in the form norm
$X\to X^*$. If \eqref{eq:r9rel-uniformcoercivity} failed, translate
a sequence of optimizers to obtain $Q_n\to Q$ strongly in $X$, and
choose real $b_n\perp Q_n$ with $\|b_n\|_X=1$ and
$\langle b_n,L_{Q_n}b_n\rangle\to0$.
The projection of $b_n$ onto $Q$ tends to zero. Form convergence and
the strict coercivity for this fixed $Q$ contradict the unit norm.

For a nonzero direction $\zeta$, the identity
$\zeta\cdot\nabla Q=0$ would give
$(\zeta\cdot\xi)\widehat Q(\xi)=0$.
An $L^2$ function supported on a hyperplane is zero, contrary to
$M(Q)=M_c$. Compactness modulo translations and continuity of
differentiation from $X$ to $X^*$ make this lower bound uniform over
$Q$ and over the unit sphere of directions.
Finally, real and purely imaginary functions are orthogonal for
the real-even multiplier defining the $X^*$ inner product.
The norms $\|Q\|_{X^*}$ have a uniform positive lower bound by
optimizer compactness, which proves
\eqref{eq:r9rel-parameterinjective}.
\end{proof}

\begin{lemma}[Uniform proximity and phase choice]\label{lem:r9rel-proximity}
There is a modulus $\varepsilon(\delta)\to0$ as $\delta\downarrow0$
such that
\[
 A(U)=1,\qquad E(U)<0,\qquad M_c<M(U)<M_c+\delta
\]
imply that, after a translation and a constant phase, some
$Q\in\mathcal Q$ satisfies
\[
 \|U-Q\|_X\le\varepsilon(\delta),\qquad
 \langle\operatorname{Im}U,Q\rangle=0.
\]
\end{lemma}
\begin{proof}
Failure of the first assertion would contradict normalized
near-optimizer compactness. A further phase tending to zero makes
$\int UQ$ real and positive; this is possible because the integral
tends to $M_c$. It changes the proximity modulus only by a fixed
factor and gives the asserted orthogonality.
\end{proof}

\subsection{A quantitative residual theorem}

\begin{theorem}[Positive inertia with a residual error]
\label{thm:r9rel-residual}
There exist $\delta,c,C>0$, depending only on $N,m,a$, such that
the following holds. Let $U\in X$ satisfy
\[
 M_c<M(U)<M_c+\delta,\qquad E(U)<0.
\]
For arbitrary $\omega\in\mathbb R$ and $v\in\mathbb R^N$, define
\[
 F=a(D)U+\omega U+iv\cdot\nabla U-|U|^\alpha U\in X^*.
\]
Then
\begin{equation}\label{eq:r9rel-residualinertia}
 v\cdot P(U)\ge cA(U)^{2/m-1}|v|^2
 -C\int\frac{|\widehat F(\xi)|^2}{A(U)+a(\xi)}\,d\xi.
\end{equation}
The constants are independent of frequency and velocity.
No radial symmetry or spatial moment is required.

If, in addition, the scalar dilation defect satisfies
\begin{equation}\label{eq:r9rel-defecthyp}
 D:=mE(U)-\tfrac12v\cdot P(U)
 \ge-\theta m|E(U)|,
 \qquad 0\le\theta<1,
\end{equation}
then, for some $c_\theta>0$ depending only on the model and $\theta$,
\begin{equation}\label{eq:r9rel-residuallower}
 \int\frac{|\widehat F(\xi)|^2}{A(U)+a(\xi)}\,d\xi
 \ge c_\theta\bigl(|E(U)|+A(U)^{2/m-1}|v|^2\bigr).
\end{equation}
\end{theorem}
\begin{proof}
First normalize $A(U)=1$, and put $\rho=\|F\|_{X^*}$.
Use Lemma~\ref{lem:r9rel-proximity} to arrange proximity to a real
$Q\in\mathcal Q$ and the stated phase orthogonality. All quantities
in the conclusion are invariant under these translations and phases.

Write $G(U)=a(D)U-|U|^\alpha U$. Its $X^*$ norm is bounded uniformly
in the normalized mass window. Replacing $Q$ by $U$ in
\eqref{eq:r9rel-parameterinjective} costs at most
$C\varepsilon(\delta)(|\omega|+|v|)$.
For sufficiently small $\delta$, absorption in
$F=G(U)+\omega U+iv\cdot\nabla U$ therefore gives
\begin{equation}\label{eq:r9rel-parameterbound}
 |\omega|+|v|\le C(1+\rho).
\end{equation}
Subtracting $G(Q)=-\omega_*Q$, and using the local Lipschitz
continuity of $G$, now gives
\begin{equation}\label{eq:r9rel-parametersmall}
 |\omega-\omega_*|+|v|
 \le C\{\rho+\varepsilon(\delta)(1+\rho)\}.
\end{equation}

Fix a sufficiently small $\rho_0>0$, and then shrink $\delta$.
For $\rho\le\rho_0$, the form
$L_U=a(D)+\omega-|U|^\alpha$ is uniformly close to $L_Q$.
Consequently,
\begin{equation}\label{eq:r9rel-perturbedcoercivity}
 \langle b,L_Ub\rangle\ge c_1\|b\|_X^2
 \quad\text{for real }b\perp Q,
\end{equation}
with $c_1>0$ independent of $U,\omega,v$ in this regime.
Write $U=h+ib$. Its imaginary residual equation is
\[
 \operatorname{Im}F=L_Ub+v\cdot\nabla h.
\]
The identity $P(U)=-2\langle b,\nabla h\rangle$ is valid by $m>1$.
Thus
\begin{align}
 \tfrac12v\cdot P(U)
 &=\langle b,L_Ub\rangle-\langle\operatorname{Im}F,b\rangle
 \notag\\
 &\ge\tfrac12c_1\|b\|_X^2-C\rho^2.
 \label{eq:r9rel-phasepairing}
\end{align}
The imaginary equation also gives, by
\eqref{eq:r9rel-derivativeinjective},
\[
 \kappa|v|\le\|v\cdot\nabla Q\|_{X^*}
 \le\rho+C\|b\|_X+C\varepsilon(\delta)|v|.
\]
Here $L_U:X\to X^*$ has uniformly bounded norm.
Absorbing the last term yields
$|v|\le C(\rho+\|b\|_X)$. Combining this with
\eqref{eq:r9rel-phasepairing} proves
\begin{equation}\label{eq:r9rel-normalized}
 v\cdot P(U)\ge c|v|^2-C\rho^2
\end{equation}
when $\rho\le\rho_0$.

If $\rho>\rho_0$, normalized mass and kinetic energy give
$|P(U)|\le C$, while \eqref{eq:r9rel-parameterbound} gives
$|v|\le C(1+\rho)$. Hence
\[
 c|v|^2-v\cdot P(U)\le C(1+\rho)^2
 \le C_{\rho_0}\rho^2.
\]
After enlarging $C$, \eqref{eq:r9rel-normalized} holds for every
residual size.

For general $A=A(U)>0$, set $\lambda=A^{-1/m}$ and
\[
 U_\lambda(x)=\lambda^{N/2}U(\lambda x),\qquad
 \omega_\lambda=\lambda^m\omega,\qquad
 v_\lambda=\lambda^{m-1}v.
\]
Then $A(U_\lambda)=1$,
$v_\lambda\cdot P(U_\lambda)=A^{-1}v\cdot P(U)$, and the
residual scales as
$F_\lambda(x)=\lambda^{N/2+m}F(\lambda x)$.
A change of Fourier variables gives exactly
\[
 \|F_\lambda\|_{X^*}^2
 =A^{-1}\int\frac{|\widehat F(\xi)|^2}{A+a(\xi)}\,d\xi.
\]
Rescaling \eqref{eq:r9rel-normalized} proves
\eqref{eq:r9rel-residualinertia}. Finally,
\eqref{eq:r9rel-defecthyp} implies
\[
 v\cdot P(U)=2mE(U)-2D
 \le-2(1-\theta)m|E(U)|.
\]
Substitution into \eqref{eq:r9rel-residualinertia} proves
\eqref{eq:r9rel-residuallower}.
\end{proof}

\begin{remark}
The defect $D$ is defined without a spatially weighted pairing of
$F$. When that pairing is justified,
$D=\operatorname{Re}\langle F,(x\cdot\nabla+N/2)U\rangle$;
the same conclusion holds through the corresponding cutoff limit.
For an exact traveling profile the Pohozaev identity gives $D=0$.
The residual theorem therefore refines the profile exclusion by
quantifying how far an approximate relative equilibrium with small
dilation defect must remain from satisfying its profile equation.
\end{remark}

\subsection{The positive response tensor}

Let $L_Q^{-1}$ denote the coercive form inverse on the real
$L^2$-orthogonal complement of $Q$. Since
$\langle\partial_jQ,Q\rangle=0$, the matrix
\begin{equation}\label{eq:r9rel-tensor}
 T_{Q,ij}=\langle\partial_iQ,L_Q^{-1}\partial_jQ\rangle
\end{equation}
is well defined.

\begin{proposition}[Uniform effective inertia]\label{prop:r9rel-tensor}
The matrices $T_Q$ are symmetric positive definite, with eigenvalues
bounded above and below by positive model-dependent constants,
uniformly for $Q\in\mathcal Q$.
Suppose
\[
 A(U_n)=1,\quad E(U_n)<0,\quad M(U_n)\downarrow M_c,
 \qquad \rho_n=\|F_n\|_{X^*}\longrightarrow0,
\]
where $F_n$ uses arbitrary parameters $(\omega_n,v_n)$.
For the optimizers $Q_n$ supplied by
Lemma~\ref{lem:r9rel-proximity},
\begin{equation}\label{eq:r9rel-tensorresponse}
 P(U_n)=2T_{Q_n}v_n+o(|v_n|)+O(\rho_n),
\end{equation}
with a uniform constant in the $O$ term.
\end{proposition}
\begin{proof}
The $L^2$ projection onto $Q^\perp$ is uniformly bounded on $X$
and its dual. Thus the restricted dual norm is uniformly equivalent
to the full norm for functionals annihilating $Q$.
For $d=\zeta\cdot\nabla Q$, uniform coercivity and boundedness of
$L_Q$ give
\[
 c\|d\|_{X^*}^2\le\langle d,L_Q^{-1}d\rangle
 \le C\|d\|_{X^*}^2.
\]
Lemma~\ref{lem:r9rel-uniform} proves both eigenvalue bounds.

Arrange $U_n=h_n+ib_n$ as in the residual proof.
From \eqref{eq:r9rel-perturbedcoercivity} and the exact momentum
pairing one also obtains
\[
 c_1\|b_n\|_X^2
 \le\tfrac12v_n\cdot P(U_n)+\rho_n\|b_n\|_X
 \le C|v_n|\|b_n\|_X+\rho_n\|b_n\|_X.
\]
Hence $\|b_n\|_X\le C(|v_n|+\rho_n)$.
Using the imaginary residual equation, form convergence, and
\eqref{eq:r9rel-parametersmall},
\[
 L_{Q_n}b_n=-v_n\cdot\nabla Q_n
 +o_{X^*}(|v_n|)+O_{X^*}(\rho_n).
\]
Project onto $Q_n^\perp$ and invert $L_{Q_n}$.
Pairing the resulting identity with $-2\nabla h_n$, and using
$\|h_n-Q_n\|_X\to0$, proves
\eqref{eq:r9rel-tensorresponse}.
\end{proof}

In particular, a residual satisfying $\rho_n=o(|v_n|)$ forces
strictly positive velocity--momentum pairing at leading order.
This conclusion uses neither uniqueness of the optimizer nor
nondegeneracy of the real-part linearized operator.

\subsection{Relative-periodic solutions and shape speed}

The following statements impose explicit regularity sufficient for
the differentiated cutoff virial identity. They do not assume that
every energy-space solution has this additional regularity.

\begin{theorem}[Nonvanishing shape speed]\label{thm:r9rel-periodic}
Suppose a solution of
\[
 iu_t=a(D)u-|u|^\alpha u
\]
has the relative-periodic representation
\[
 u(t,x)=e^{i\omega t}V(t,x-vt),\qquad V(t+T)=V(t),
\]
where $T>0$ and
\begin{equation}\label{eq:r9rel-periodicregularity}
 V\in C_{\mathrm{per}}([0,T];H^{m+1})
 \cap C^1_{\mathrm{per}}([0,T];H^{m/2}).
\end{equation}
Assume that the conserved energy is negative and that
$M_c<M<M_c+\delta$, with $\delta$ as in
Theorem~\ref{thm:r9rel-residual}. Then
\begin{equation}\label{eq:r9rel-periodicpohozaev}
 v\cdot P=2mE<0.
\end{equation}
At every time and for every real number $\nu$,
\begin{equation}\label{eq:r9rel-shapespeed}
 \int\frac{|\widehat{\partial_tV-i\nu V}(t,\xi)|^2}
 {A(V(t))+a(\xi)}\,d\xi
 \ge c\bigl(|E|+A(V(t))^{2/m-1}|v|^2\bigr),
\end{equation}
where $c>0$ depends only on the model.
\end{theorem}
\begin{proof}
The equation in this frame is
\[
 iV_t=a(D)V+\omega V+iv\cdot\nabla V-|V|^\alpha V.
\]
Choose a real $\chi\in C_c^\infty$ equal to one near the origin,
and set
\[
 f_R(x)=x\chi(x/R),\qquad
 L_R=f_R\cdot\nabla+\tfrac12\operatorname{div}f_R,
 \qquad J_R(V)=\operatorname{Im}\int\overline V L_RV.
\]
The operator $L_R$ is skew-symmetric. Differentiation gives
\begin{align}
 J_R'(t)
 &=\langle V,[a(D),L_R]V\rangle
 +\langle V,[iv\cdot\nabla,L_R]V\rangle\notag\\
 &\quad-2\left(\tfrac12-\tfrac1p\right)
 \int\operatorname{div}f_R\,|V|^p.
 \label{eq:r9rel-cutoffvirial}
\end{align}
The frequency term vanishes by skew symmetry.
Assumption~\eqref{eq:r9rel-periodicregularity} justifies the
differentiation and all pairings. In particular,
$|V|^\alpha V\in L^2$ by Sobolev embedding from $H^{m+1}$,
and $\nabla V\in L^2$, so the nonlinear integration by parts is valid.

For completeness, the kinetic commutator has Fourier kernel,
up to the fixed normalization factor,
\[
 \frac i2\bigl(a(\xi)-a(\eta)\bigr)
 (\xi+\eta)\cdot\widehat f_R(\xi-\eta).
\]
Homogeneity gives
$|a(\xi)-a(\eta)|\le C|\xi-\eta|(|\xi|+|\eta|)^{m-1}$.
After division by the two energy weights, the kernel is bounded
by $C|k|\langle k\rangle^{m/2}|\widehat f_R(k)|$,
where $k=\xi-\eta$. Since
$\widehat f_R(k)=R^{N+1}\widehat f(Rk)$, this majorant has
uniformly bounded $L^1$ norm for $R\ge1$.
Young's inequality gives a uniform bound $X\to X^*$.
For Schwartz functions, Fourier integration by parts and
$\xi\cdot\nabla a=ma$ give the limiting form $mA$.
Density and the uniform bound extend the limit to every $V\in X$.
The identical degree-one argument for $iv\cdot\nabla$ gives
the limit $-v\cdot P$. Finally, dominated convergence gives the
nonlinear limit $-2m\|V\|_p^p/p$.
Thus the right side of \eqref{eq:r9rel-cutoffvirial} tends to
\[
 mA(V)-v\cdot P-\frac{2m}{p}\|V\|_p^p
 =2mE-v\cdot P.
\]
Uniform form bounds and boundedness of the trajectory on the period
permit dominated convergence in time. Integrate first over
$[0,T]$. Periodicity gives $J_R(T)-J_R(0)=0$ for every $R$;
then let $R\to\infty$. This proves
\eqref{eq:r9rel-periodicpohozaev} without a spatial moment.

At a fixed time apply Theorem~\ref{thm:r9rel-residual} to $V(t)$
with frequency $\omega+\nu$ and velocity $v$.
Its residual is
\[
 F=iV_t+\nu V=i(V_t-i\nu V).
\]
Equation~\eqref{eq:r9rel-periodicpohozaev} gives $D=0$.
The lower bound \eqref{eq:r9rel-residuallower}, with $\theta=0$,
is precisely \eqref{eq:r9rel-shapespeed}.
\end{proof}

The result also applies to a profile periodic up to a constant
phase: a linear phase adjustment makes the representation periodic,
and the cutoff functional is phase invariant.

\begin{corollary}[Moving-center cost over a cycle]
\label{cor:r9rel-movingcenter}
Under the same mass and energy hypotheses, suppose instead that
\[
 W(t,x)=e^{i\gamma(t)}u(t,x+c(t)),\qquad W(t+T)=W(t),
\]
where $c,\gamma$ are $C^1$ and $W$ has the regularity
\eqref{eq:r9rel-periodicregularity}.
Put $d=c(T)-c(0)$. Then
\begin{equation}\label{eq:r9rel-netdisplacement}
 P\cdot d=2mET,
\end{equation}
and
\begin{align}
 &\int_0^T\int
 \frac{|\widehat{W_t}(t,\xi)|^2}{A(W(t))+a(\xi)}\,d\xi\,dt
 \notag\\[-2pt]
 &\hspace{15mm}\ge
 c\left\{|E|T+
 \int_0^T A(W(t))^{2/m-1}|c'(t)|^2\,dt\right\}.
 \label{eq:r9rel-cyclecost}
\end{align}
The second term in braces is at least
\begin{equation}\label{eq:r9rel-displacementcost}
 \frac{|d|^2}{\displaystyle\int_0^T A(W(t))^{1-2/m}\,dt}.
\end{equation}
\end{corollary}
\begin{proof}
The equation for $W$ is
\[
 iW_t=a(D)W-\gamma'(t)W+ic'(t)\cdot\nabla W-|W|^\alpha W.
\]
The same cutoff virial calculation, integrated over the period,
gives \eqref{eq:r9rel-netdisplacement}.
At each time, Theorem~\ref{thm:r9rel-residual} applies with
frequency $-\gamma'(t)$, velocity $c'(t)$, and residual $iW_t$.
Integrating \eqref{eq:r9rel-residualinertia}, and substituting
\eqref{eq:r9rel-netdisplacement}, proves
\eqref{eq:r9rel-cyclecost}. Weighted Cauchy--Schwarz gives
\eqref{eq:r9rel-displacementcost}.
\end{proof}

These inequalities give necessary quantitative restrictions on
near-threshold negative-energy coherent dynamics. They exclude
vanishing shape speed in the stated relative-periodic class, even
after a time-dependent choice of center and phase. They neither
exclude all breathers nor establish a mass gap for arbitrary bounded
global solutions. The left side of \eqref{eq:r9rel-cyclecost} measures
evolution in a dual energy norm; it is not a physical dissipation,
and no loss of conserved energy is asserted.


\section{What remains open in the dynamical continuation}
For pure BNLS, the specific unresolved step in this continuation is whether a uniform mass interval above $M_c$ excludes every forward-global, uniformly $H^2$-bounded negative-energy solution with zero momentum. The profile theorem only covers relative equilibria. The relative-periodic estimates impose restrictions on a larger class, while the virial counterexample excludes one uniform pointwise closure. None provides the missing classification or compactness statement for arbitrary bounded trajectories.

In particular, the following two inferences are invalid: that absence of traveling profiles excludes arbitrary global solutions; or that a small near-optimizer tail can always be absorbed into the negative energy. The first requires a theorem about the actual dynamics. The second fails for the explicit localized interaction functional analyzed below, even for smooth zero-momentum data arbitrarily close to an optimizer. Any successful continuation through that method must exploit information from the evolution beyond a pointwise estimate on all such functions.

\section{Relation to prior work}
The complex-to-real optimizer argument has precedent in Fern\'andez--Jeanjean--Mandel--Mari\c{s}~\cite{FJMM}, Proposition 3.8. Related variational phase-kernel arguments appear in Iyer--Parker--Stefanov~\cite{IPS}, Section 3.3.2, for a one-dimensional fourth-order problem. Appendix~\ref{sec:foundations} supplies the general homogeneous-symbol algebra needed here. Concentration compactness, sharp interpolation, phase rotation and constrained minimization are established methods.

Chen--Wang~\cite{CW}, Theorem 1.4, study biharmonic normalized stationary solutions with a critical focusing power and a lower defocusing power. Their nonexistence statement at or below the critical mass and their constrained-energy statements do not assert a uniform strict interval above that mass for all negative-energy traveling profiles. Their parameter $a$ is an $L^2$ norm, whereas mass in this manuscript is its square.

The higher-order comparison in Appendix~\ref{sec:superorder} has direct precedents. Bonheure--Cast\'eras--Moreira dos Santos--Nascimento~\cite{BCMN} prove normalized-wave existence for mixed second- and fourth-order dispersion, including the lower-order mass-critical exponent. Luo--Hajaiej~\cite{LH} obtain the sharp mass-threshold alternative for mixed fractional Laplacians, including the lower-order critical endpoint. Thus the higher-order stabilization mechanism and its variational proof are not presented as newly discovered physical principles. The appendix gives the arbitrary-order version and bounds needed for the present comparison.

Pausader--Shao~\cite{PS}, Section 6, already connect center displacement, momentum, energy and interaction virial identities for compact trajectories. The recent Inui--Komada--Nishimura theorem~\cite{IKN} proves nonradial scattering in the range $p\ge1+8/(d-1)$, with its other hypotheses, and records the averaged center-motion identity. The mass-critical exponent $1+8/d$ considered here is outside that scattering theorem. The center-motion identity itself is not a novelty claim of this manuscript; the residual estimate supplies an additional quantitative cost of shape motion.

The contribution to be assessed is the combined uniform exclusion theorem and its shifted-frequency proof, the explicit same-order small-perturbation boundary, and the quantitative residual conclusion in their stated generality. The targeted primary-source comparison is not an exhaustive proof of priority.

\appendix
\section{Variational foundations}\label{sec:foundations}

For completeness, this appendix supplies the two variational ingredients used above. The notation is that of Section~\ref{sec:cone}; all spaces may be complex unless real functions are specified. These arguments require neither a maximum principle nor uniqueness of an optimizer.


\begin{lemma}[Compactness of normalized near-optimizers]
\label{foundation:compact}
The sharp inequality \eqref{cone:gn} has an optimizer.
If $M(f_n)\to M_{c,a}$, $A_a(f_n)\to1$, and $\|f_n\|_p^p\to p/2$,
then after translations and passage to a subsequence,
$f_n$ converges strongly in $H^{m/2}$ to a normalized optimizer. After a constant phase, that optimizer is real. The family of normalized optimizers is compact modulo translations and constant phases.
\end{lemma}
\begin{proof}
The hypotheses make the sequence a normalized maximizing sequence. In particular, the same conclusion holds if $A_a(f_n)=1$ and $E_a(f_n)<0$, since the sharp inequality then forces $\|f_n\|_p^p\to p/2$.
Local Sobolev nonvanishing gives translations for which a subsequence
has a nonzero weak $H^{m/2}$ limit $Q$. Local Rellich compactness
also gives almost-everywhere convergence after a subsequence.
Write
\[
 m_1=M(Q)>0,\quad a_1=A_a(Q)>0,\quad
 m_2=M_{c,a}-m_1,\quad a_2=1-a_1.
\]
Hilbert norm splitting, the Brezis--Lieb lemma, and
\eqref{cone:gn} give
\[
 M_{c,a}^r\le m_1^ra_1+m_2^ra_2,\qquad a_1+a_2=1.
\]
If $0<m_1<M_{c,a}$, both coefficients $m_j^r$ are strictly below
$M_{c,a}^r$, which is impossible. Hence $m_1=M_{c,a}$, and the
same inequality forces $a_1=1$. Strong energy-space convergence
and equality in the sharp inequality follow.
To prove attainment, normalize an arbitrary maximizing sequence
by amplitude and dilation so that $M(f_n)=M_{c,a}$ and
$A_a(f_n)=1$. Its nonlinear norms converge to $p/2$, and the
identical splitting argument applies; negative energy is not
needed once that convergence is known. The same proof applies to a sequence of normalized optimizers. The constant-phase assertion follows from the next lemma.
\end{proof}

\begin{lemma}[Constant phase of optimizers]\label{foundation:constant-phase}
Every complex optimizer is a constant phase times a real function.
\end{lemma}
\begin{proof}
Write $f=A+iB$, where $A,B$ are real. Evenness of $a$ gives
$A_a(f)=A_a(A)+A_a(B)$. With
$C=(p/2)M_{c,a}^{-r}$, Minkowski, the sharp inequality on each real
component, and two-term H\"older yield
\begin{align*}
 \|f\|_p^2
 &\le\|A\|_p^2+\|B\|_p^2\\
 &\le C^{2/p}\left[
 A_a(A)^{2/p}M(A)^{1-2/p}
 +A_a(B)^{2/p}M(B)^{1-2/p}\right]\\
 &\le C^{2/p}A_a(f)^{2/p}M(f)^{1-2/p}.
\end{align*}
For an optimizer all inequalities are equalities. Strict convexity
of $L^{p/2}$ implies $B^2=cA^2$ almost everywhere, unless one
component vanishes. Repeat this equality argument after multiplying
$f$ by $e^{i\pi/4}$. If both signs of $B/A=\pm\sqrt c$ occurred
on sets of positive measure, the new squared components could
not be proportional. When $c=1$ the new components instead have
disjoint nonzero supports, again giving strict Minkowski inequality.
Thus only one sign is possible, and $f$ has constant phase.
\end{proof}

Let $Q$ be a real optimizer normalized by $M(Q)=M_{c,a}$ and
$A_a(Q)=1$. Its Euler--Lagrange equation is
\[
 a(D)Q+\omega_QQ-|Q|^{2r}Q=0,\qquad
 \omega_Q=r/M_{c,a}>0.
\]

\begin{lemma}[Strict phase coercivity]\label{foundation:phase}
The real quadratic form
\[
 L_-=a(D)+\omega_Q-|Q|^{2r}
\]
is nonnegative, has kernel $\operatorname{span}_{\mathbb R}\{Q\}$,
and obeys
\[
 \langle b,L_-b\rangle\ge c_Q\|b\|_{H^{m/2}}^2
 \quad\text{if }b\text{ is real and }\langle b,Q\rangle=0,
\]
for some $c_Q>0$.
\end{lemma}
\begin{proof}
Retain $A=A_a(Q)$ and $M=M(Q)$ in the algebra, so that
$\omega_Q=rA/M$. For real $b\in H^{m/2}$ set
$B=M(b)$ and $C_1=A_a(b)$.
The sharp inequality applied to $Q+i\sqrt t\,b$ gives
\begin{equation}\label{foundation:kernel-comparison}
 F(t):=\int(Q^2+tb^2)^{r+1}
 \le(r+1)(1+tB/M)^r(A+tC_1),\qquad t\ge0.
\end{equation}
The right derivatives at zero prove $L_-\ge0$.
If $b$ belongs to its form kernel, the first derivatives agree, and
\[
 \int |Q|^{2r}b^2=C_1+rAB/M.
\]
The pointwise second-order remainder for $F$ is nonnegative.
Comparison of second-order remainders in
\eqref{foundation:kernel-comparison} and Fatou's lemma yield
\[
 D:=\int |Q|^{2r-2}b^4
 \le(r-1)AB^2/M^2+2BC_1/M.
\]
On the zero set of $Q$, the integrand is interpreted through
the remainder: it is infinite if $r<1$ and $b\ne0$, equals $b^4$
if $r=1$, and equals zero if $r>1$.
Cauchy--Schwarz gives
\[
 D\ge\frac{(C_1+rAB/M)^2}{(r+1)A}.
\]
Combining the inequalities yields
$(C_1-AB/M)^2\le0$. Consequently all inequalities are equalities.
Equality in Cauchy--Schwarz implies
$b^2=(B/M)Q^2$ on $\{Q\ne0\}$. Its integral already equals $B$,
so $b=0$ on the remaining zero set as well.
The form kernel is a linear space containing $Q$. Apply the same
conclusion to $b+Q$ and subtract the identity for $b^2$.
It follows that $b$ is one constant multiple of $Q$.

Finally, multiplication by $|Q|^{2r}$ is compact from $H^{m/2}$
to $H^{-m/2}$. Indeed $|Q|^{2r}\in L^{p/(2r)}$,
H\"older with two $L^p$ factors makes the tail small, and the
embedding into local $L^p$ is compact because $p$ lies strictly
below the critical Sobolev exponent. Thus $L_-$ is a compact
form perturbation of the coercive form $a(D)+\omega_Q$.
Nonnegativity and its one-dimensional kernel give the asserted
coercivity on the $L^2$ orthogonal complement of $Q$.
\end{proof}


\section{A near-threshold obstruction to a pointwise interaction-virial sign}
\label{sec:r9-virial-obstruction}

This section concerns the pure mass-critical BNLS
\[
 iu_t=\Delta^2u-|u|^{8/N}u,\qquad N\ge1.
\]
Write
\[
 p=2+8/N,\quad M(f)=\|f\|_2^2,\quad A(f)=\|\Delta f\|_2^2,
 \quad E(f)=\frac12A(f)-\frac1p\|f\|_p^p,
\]
\[
 P(f)=\operatorname{Im}\int\overline f\nabla f,
 \qquad W(f)=\operatorname{Im}\int\Delta f\nabla\overline f.
\]
Let $M_c$ be the sharp Gagliardo--Nirenberg critical mass. Fix the
interaction cutoff $\phi_R$ which equals $|z|^2/2$ for $|z|\le4R$,
is constant for $|z|\ge8R$, and satisfies $D^2\phi_R\le I$. Define
\begin{equation}\label{eq:r9-interaction-functional}
 I_R(f)=2\iint |f(y)|^2\nabla\phi_R(x-y)\cdot
       \operatorname{Im}(\overline f(x)\nabla f(x))\,dx\,dy.
\end{equation}
For smooth compactly supported $f$, let $\mathcal D_R(f)$ denote
the initial derivative of this functional prescribed by the BNLS
vector field, equivalently the right-hand side of the exact
interaction-virial identity evaluated at $f$.

\begin{theorem}[Positive localized virial derivative arbitrarily near threshold]
\label{thm:r9-static-obstruction}
For every $\delta>0$, there exist constants
$M_*,A_*,e_*,c_*,R_*>0$ and functions
$f_R\in C_c^\infty(\mathbb R^N;\mathbb C)$, $R\ge R_*$, such that
\begin{align}
 M_c<M(f_R)&=M_*<M_c+\delta,&
 A(f_R)&=A_*,\label{eq:r9-static-invariants}\\
 P(f_R)&=0,&E(f_R)&=-e_*<0,\notag
\end{align}
and
\begin{equation}\label{eq:r9-static-positive}
 \mathcal D_R(f_R)\ge c_*>0.
\end{equation}
All constants are independent of $R$. Moreover, one may choose
these families as $\delta\downarrow0$ so that, for a fixed real
normalized optimizer $Q$,
\[
 \sup_{R\ge R_*(\delta)}\|f_R-Q\|_{H^2}\longrightarrow0.
\]
Consequently no estimate
\[
 \mathcal D_R(f)\le c\,M(f)E(f)+o_{R\to\infty}(1),\qquad c>0,
\]
can hold uniformly over zero-momentum negative-energy data of
sufficiently near-critical mass and bounded kinetic energy.
\end{theorem}

\begin{lemma}[Independent momentum and flux directions]
\label{lem:r9-independent-currents}
For every nonzero real $q\in C_c^\infty(\mathbb R^N)$, the linear
map
\[
 b\longmapsto \bigl(P(q+ib),W(q+ib)\bigr),
 \qquad b\in C_c^\infty(\mathbb R^N;\mathbb R),
\]
is onto $\mathbb R^N\times\mathbb R^N$.
\end{lemma}
\begin{proof}
Integration by parts gives
\begin{equation}\label{eq:r9-current-linearization}
 P(q+ib)=-2\int b\nabla q,\qquad
 W(q+ib)=2\int b\nabla\Delta q.
\end{equation}
The $2N$ real functions $\partial_jq,\partial_j\Delta q$ are
linearly independent in $L^2$. Indeed, a dependence would imply
\[
 \bigl[(a\cdot\xi)-(b_0\cdot\xi)|\xi|^2\bigr]\widehat q(\xi)=0
 \quad\text{a.e.}
\]
for some $a,b_0\in\mathbb R^N$. A nonzero polynomial has a
Lebesgue-null zero set, whereas a nonzero $L^2$ function cannot
be supported on a null set. Thus the polynomial vanishes
identically, forcing $a=b_0=0$. The Gram matrix of the $2N$
functions is therefore invertible. Taking $b$ in their span
realizes every pair in \eqref{eq:r9-current-linearization}.
\end{proof}

\begin{lemma}[Separated-packet identity]
\label{lem:r9-separated-virial}
Suppose $F,G\in C_c^\infty$ have support diameters below $4R$, and
the distance between their supports is above $8R$. Then
\begin{equation}\label{eq:r9-separated-identity}
 \mathcal D_R(F+G)=\mathcal V(F)+\mathcal V(G),\qquad
 \mathcal V(h)=16M(h)E(h)-8P(h)\cdot W(h).
\end{equation}
\end{lemma}
\begin{proof}
For intrapacket pairs the cutoff is exactly quadratic; for
interpacket pairs it is constant. Substitution into the exact
interaction identity gives \eqref{eq:r9-separated-identity}.
Equivalently, the quadratic interaction associated with one
packet is
\[
 2\left[M(h)\operatorname{Im}\int x\cdot\overline h\nabla h
       -\left(\int x|h|^2\right)\cdot P(h)\right].
\]
With $J(h)=\operatorname{Im}\int x\cdot\overline h\nabla h$,
one has $J'=8E(h)$, while the mass first moment has derivative
$4W(h)$. Mass and momentum are conserved, so the displayed
interaction has derivative $2[8ME-4W\cdot P]$, as asserted. This is an
initial-derivative calculation; no finite-speed propagation or
persistence of compact support is asserted.
\end{proof}

\begin{proof}[Proof of Theorem~\ref{thm:r9-static-obstruction}]
Fix a real optimizer $Q$ with $M(Q)=M_c$ and $E(Q)=0$.
Approximate it in $H^2$ by a nonzero real compactly supported
smooth function $q_0$. Multiplication by
\[
 \left(\frac{pA(q_0)}{2\|q_0\|_p^p}\right)^{1/(p-2)}
\]
gives a real $q\in C_c^\infty$ satisfying
\begin{equation}\label{eq:r9-zero-energy-core}
 E(q)=0,\qquad M_0:=M(q)<M_c+\delta/4.
\end{equation}
The multiplier tends to one as $q_0\to Q$.

By Lemma~\ref{lem:r9-independent-currents}, there are fixed real
$b_1,b_2\in C_c^\infty$ such that
\[
 P(q+ib_1)=0,\quad W(q+ib_1)=e_1,\qquad
 P(q+ib_2)=-e_1,\quad W(q+ib_2)=0.
\]
Set $b_\varepsilon=\varepsilon b_1+\varepsilon^2b_2$. Then
\begin{equation}\label{eq:r9-prescribed-core-currents}
 P(q+ib_\varepsilon)=-\varepsilon^2e_1,\qquad
 W(q+ib_\varepsilon)=\varepsilon e_1,\qquad
 \|b_\varepsilon\|_{H^2}=O(\varepsilon).
\end{equation}
Since
\[
 \left.\partial_aE(aq)\right|_{a=1}
 =A(q)-\|q\|_p^p=-\frac4N A(q)<0,
\]
the implicit-function theorem supplies $a_\varepsilon=1+O(\varepsilon^2)$
for which
\begin{equation}\label{eq:r9-core-negative-energy}
 F_\varepsilon=a_\varepsilon q+ib_\varepsilon,
 \qquad E(F_\varepsilon)=-\frac{\varepsilon^3}{4M_0}.
\end{equation}
Here the energy is $C^2$ on this finite-dimensional family because
$p>2$, and its change under the imaginary perturbation starts at
second order. Write $M_F=M(F_\varepsilon)$. Exactly,
\[
 M_F=a_\varepsilon^2M_0+\|b_\varepsilon\|_2^2,
 \qquad P_F=-a_\varepsilon\varepsilon^2e_1,
 \qquad W_F=a_\varepsilon\varepsilon e_1.
\]
Consequently
\begin{equation}\label{eq:r9-core-positive-virial}
 \mathcal V(F_\varepsilon)
 =\left(8a_\varepsilon^2-\frac{4M_F}{M_0}\right)\varepsilon^3
 =(4+O(\varepsilon^2))\varepsilon^3>2\varepsilon^3
\end{equation}
for all sufficiently small positive $\varepsilon$.

Fix $\mathfrak m\in(0,\min\{M_c,\delta/4\})$, and choose a real
radial $\chi\in C_c^\infty(\mathbb R^N)$ with $\|\chi\|_2=1$.
Define a second packet by
\begin{equation}\label{eq:r9-balancing-packet}
 G(x)=\sqrt{\mathfrak m}\,L^{-N/2}\chi(x/L)e^{ikx_1},
 \qquad k=\frac{a_\varepsilon\varepsilon^2}{\mathfrak m}.
\end{equation}
Its mass is $\mathfrak m$ and its momentum is exactly $-P_F$.
Put
\[
 B=\|\nabla\chi\|_2^2,\quad C=\|\Delta\chi\|_2^2,
 \quad D=\|\chi\|_p^p,\quad d=1+2/N.
\]
Fourier expansion and radial symmetry give
\begin{align}
 A(G)&=\mathfrak m\bigl[k^4+2d k^2L^{-2}B+L^{-4}C\bigr],
 \label{eq:r9-radiation-formulas}\\
 W(G)&=\mathfrak m\bigl[k^3+d kL^{-2}B\bigr]e_1,
 &\|G\|_p^p&=\mathfrak m^{p/2}L^{-4}D.\notag
\end{align}
In particular, the plane-wave terms cancel exactly in the
self-interaction derivative:
\begin{equation}\label{eq:r9-radiation-positive-virial}
 \mathcal V(G)=8\mathfrak m^2
 \left[d k^2L^{-2}B+
 L^{-4}\left(C-\frac2p\mathfrak m^{4/N}D\right)\right]>0.
\end{equation}
The last inequality follows from the sharp GN inequality and
$\mathfrak m<M_c$.
For fixed $\varepsilon$, as $L\to\infty$,
\[
 E(G)\longrightarrow\frac{\mathfrak m k^4}{2}
 =\frac{a_\varepsilon^4\varepsilon^8}{2\mathfrak m^3}.
\]
Choose $\varepsilon>0$ small enough that this limit is less than
$\varepsilon^3/(16M_0)$ and that $M_F<M_c+\delta/2$.
Then choose $L$ sufficiently large that
\begin{equation}\label{eq:r9-radiation-energy-budget}
 0<E(G)<\frac{\varepsilon^3}{8M_0}.
\end{equation}

Fix these packets. For sufficiently large $R$, set
\[
 f_R(x)=F_\varepsilon(x)+G(x-20Re_1).
\]
Their support diameters are below $4R$ and their mutual support
distance exceeds $8R$. All invariants in
\eqref{eq:r9-static-invariants} are exactly independent of $R$
by disjointness of the supports and locality of the pure BNLS
energy. Their momentum is zero, and
\[
 E(f_R)=E(F_\varepsilon)+E(G)<-\frac{\varepsilon^3}{8M_0}<0,
 \qquad M(f_R)<M_c+\delta.
\]
Negative energy forces $M(f_R)>M_c$. Lemma~\ref{lem:r9-separated-virial}
and \eqref{eq:r9-core-positive-virial}--\eqref{eq:r9-radiation-positive-virial}
prove \eqref{eq:r9-static-positive} with $c_*=2\varepsilon^3$.
Finally, as $\delta\downarrow0$, choose $q\to Q$,
$\mathfrak m\to0$, and then successively smaller $\varepsilon$
and larger $L$. The core tends to $Q$ in $H^2$, and
\eqref{eq:r9-radiation-formulas} makes $\|G\|_{H^2}\to0$.
This norm is translation invariant, proving uniform closeness
in $R$.
\end{proof}

\begin{remark}[The exact signed remainder]
\label{rem:r9-hessian-deficit}
The obstruction persists when the full favorable Hessian deficit
is retained. Let $\mathsf D_R=I-D^2\phi_R\ge0$, and set
$\rho=|f|^2$, $j=\operatorname{Im}(\overline f\nabla f)$,
$w=\operatorname{Im}(\Delta f\nabla\overline f)$. Define
\begin{align*}
 \mathcal K_R&=\iint\rho(y)(\mathsf D_R)_{jk}(x-y)
  \operatorname{Re}(\partial_{ij}\overline f\,\partial_{ik}f)(x)
  \,dx\,dy,\\
 \mathcal N_R&=\frac{2\beta}{p}\iint\rho(y)
  \operatorname{tr}\mathsf D_R(x-y)|f(x)|^p\,dx\,dy,
  \qquad\beta=8/N,\\
 \mathcal X_R&=8\iint\mathsf D_R(x-y)[j(x),w(y)]\,dx\,dy.
\end{align*}
For the support configurations above, every term containing a
higher cutoff derivative vanishes exactly, and
\[
 \mathcal D_R(f_R)=16ME-8P\cdot W
                  -8\mathcal K_R+\mathcal N_R+\mathcal X_R.
\]
Moreover,
\begin{align*}
 \mathcal K_R&=M_F A(G)+\mathfrak m A(F_\varepsilon),\\
 \mathcal N_R&=\frac{16}{p}
  \bigl(M_F\|G\|_p^p+\mathfrak m\|F_\varepsilon\|_p^p\bigr),\\
 \mathcal X_R&=8\bigl(P_F\cdot W(G)+P(G)\cdot W_F\bigr).
\end{align*}
Since $P=0$ and $\mathcal D_R(f_R)>0>16ME$, one has the strict
inequality $\mathcal N_R+\mathcal X_R>8\mathcal K_R$. Thus the
failure of a universal pointwise sign is an exact property of
the interaction identity. The theorem imposes no conclusion on
the subsequent lifespan or long-time dynamics of these data.
\end{remark}

\section{A constructive boundary: positive higher-order dispersion}
\label{sec:superorder}

A positive addition of strictly higher degree has the opposite
dilation sign and admits negative-energy stationary waves arbitrarily
close to the lower-order critical mass. This comparison lies outside
both the upper dilation condition and the original order-$m$ growth
class. The proof uses standard constrained minimization. Its
fourth/second-order and fractional-order special cases have substantial
precedent; it is included to explain the variational mechanism in
arbitrary orders, not as a counterexample to the unmodified homogeneous
equation or as an isolated sharpness test of the upper dilation condition.

Fix $N\ge1$, $k>m>1$, and $\varepsilon>0$. Set
\[
 r=\frac mN,\qquad p=2+2r,\qquad
 A_j(f)=\int |\xi|^j|\widehat f(\xi)|^2\,d\xi,
\]
and consider
\begin{equation}\label{higher:equation}
 iu_t=(|D|^m+\varepsilon|D|^k)u-|u|^{2m/N}u,
 \qquad
 E_\varepsilon(f)=\frac12A_m(f)+\frac\varepsilon2A_k(f)
                    -\frac1p\|f\|_p^p.
\end{equation}
Let $M_{c,m}$ be the sharp critical mass for the \emph{lower-order}
symbol $|\xi|^m$:
\begin{equation}\label{higher:gn}
 \|f\|_p^p\le\frac p2
       \left(\frac{M(f)}{M_{c,m}}\right)^r A_m(f).
\end{equation}
Define
\[
 I_\varepsilon(M)=\inf\{E_\varepsilon(f):
      f\in H^{k/2}(\mathbb R^N;\mathbb C),\ M(f)=M\}.
\]

\begin{theorem}[No gap above the lower-order threshold]
\label{higher:existence}
For $0<M\le M_{c,m}$, $I_\varepsilon(M)=0$ and the infimum is not
attained. For every $M>M_{c,m}$,
\[
 -\infty<I_\varepsilon(M)<0,
\]
and every minimizing sequence is precompact in $H^{k/2}$ after
translations and passage to a subsequence. There is a real minimizer
$Q_M\in L^\infty\cap H^{k+1}$ satisfying
\begin{equation}\label{higher:profile}
 (|D|^m+\varepsilon|D|^k)Q_M+\omega_MQ_M
       -|Q_M|^{2m/N}Q_M=0,
 \qquad \omega_M>0.
\end{equation}
Its energy obeys
\begin{equation}\label{higher:energy}
 E_\varepsilon(Q_M)
       =-\frac{k-m}{2m}\varepsilon A_k(Q_M)<0.
\end{equation}
Thus $e^{i\omega_Mt}Q_M$ is a global stationary solution of
\eqref{higher:equation}, with zero momentum and strictly negative
energy, for every $M>M_{c,m}$.
\end{theorem}

\begin{proof}
\emph{Coercivity and sign.}
Highest-order Gagliardo--Nirenberg interpolation gives
\begin{equation}\label{higher:coercive-gn}
 \|f\|_p^p\le C_{N,m,k}
 M(f)^{1+m/N-m/k}A_k(f)^{m/k}.
\end{equation}
Since $m/k<1$, Young's inequality implies, at fixed mass,
\begin{equation}\label{higher:coercivity}
 E_\varepsilon(f)\ge\frac\varepsilon4A_k(f)
                         -C_{N,m,k,\varepsilon,M}.
\end{equation}
The infimum is finite and minimizing sequences are bounded in
$H^{k/2}$. If $M\le M_{c,m}$, \eqref{higher:gn} gives
\[
 E_\varepsilon(f)\ge
 \frac12\left[1-(M/M_{c,m})^r\right]A_m(f)
                       +\frac\varepsilon2A_k(f)>0
\]
for every nonzero admissible $f$. On the other hand, for any smooth
datum of mass $M$, the mass-preserving dilation
$f_\lambda(x)=\lambda^{N/2}f(\lambda x)$ satisfies
\begin{equation}\label{higher:dilation}
 E_\varepsilon(f_\lambda)
       =\lambda^mE_0(f)+\frac\varepsilon2\lambda^kA_k(f)
       \longrightarrow0\quad(\lambda\downarrow0).
\end{equation}
This proves the first assertion. If $M>M_{c,m}$, multiply a
lower-order sharp optimizer by $\sqrt{M/M_{c,m}}$ to obtain negative
$E_0$. Approximation in $H^{m/2}$ by real $C_c^\infty$ functions,
followed by a mass adjustment, gives such an $f$ of exact mass $M$.
Equation~\eqref{higher:dilation} is negative for sufficiently small
$\lambda$, because $k>m$. This does not require the lower-order
optimizer itself to belong to $H^{k/2}$.

\emph{Strict binding.}
For $\theta>1$, amplitude scaling gives
\[
 E_\varepsilon(\sqrt\theta f)
 =\theta E_\varepsilon(f)
       -\frac{\theta^{p/2}-\theta}{p}\|f\|_p^p
 \le\theta E_\varepsilon(f),
\]
and also
\[
 E_\varepsilon(\sqrt\theta f)
 =\theta^{p/2}E_\varepsilon(f)
  +\frac{\theta-\theta^{p/2}}2(A_m(f)+\varepsilon A_k(f))
 \le\theta^{p/2}E_\varepsilon(f).
\]
Consequently $I_\varepsilon(M)/M$ is nonincreasing, with
\begin{equation}\label{higher:strict-amplitude}
 I_\varepsilon(\theta M)\le\theta^{p/2}I_\varepsilon(M)
               <\theta I_\varepsilon(M)
 \quad\text{if }I_\varepsilon(M)<0.
\end{equation}
For $M>M_{c,m}$ and $M=M_1+M_2$ with $M_1,M_2>0$, it follows that
\begin{equation}\label{higher:binding}
 I_\varepsilon(M)<I_\varepsilon(M_1)+I_\varepsilon(M_2).
\end{equation}
Indeed, if both pieces are subthreshold the right side is zero.
Otherwise the monotonicity of $I_\varepsilon(M)/M$, strict for a
piece with negative infimum by~\eqref{higher:strict-amplitude}, proves
\eqref{higher:binding} upon multiplying by the two masses and adding.

\emph{Compactness and attainment.}
Let $f_n$ be a minimizing sequence at $M>M_{c,m}$ and put $s=k/2$.
Local mass vanishing is impossible. To see this, a smooth unit-cube
partition and local GN interpolation give
\[
 \|\chi_j f_n\|_{2+2k/N}^{2+2k/N}
    \le C\|\chi_j f_n\|_2^{2k/N}\|\chi_j f_n\|_{H^s}^2,
 \qquad
 \sum_j\|\chi_j f_n\|_{H^s}^2\le C\|f_n\|_{H^s}^2.
\]
Thus $\sup_y\int_{B(y,1)}|f_n|^2\to0$ would imply
$\|f_n\|_{2+2k/N}\to0$, and interpolation would imply
$\|f_n\|_p\to0$ since $2<p<2+2k/N$. This contradicts the negative
limiting energy.

After translations and a subsequence, there is therefore a weak
$H^s$ limit $Q\ne0$, with local strong $L^2$ convergence and
almost-everywhere convergence. Write $w_n=f_n-Q$,
$M_1=M(Q)>0$, and $M_2=M-M_1\ge0$. Hilbert splitting of both
kinetic forms and the Brezis--Lieb lemma yield
\begin{align}
 M(f_n)&=M(Q)+M(w_n)+o(1),\notag\\
 E_\varepsilon(f_n)&=E_\varepsilon(Q)
                         +E_\varepsilon(w_n)+o(1).
 \label{higher:splitting}
\end{align}
If $M_2>0$, normalize $w_n$ to exactly that mass. Its normalization
factor tends to one and all energy terms remain bounded, so the
energy changes by $o(1)$. Equation~\eqref{higher:splitting} then
contradicts~\eqref{higher:binding}. Hence $M_1=M$ and $M(w_n)\to0$.
Equation~\eqref{higher:coercive-gn} now gives $\|w_n\|_p\to0$,
so~\eqref{higher:splitting} forces
\[
 E_\varepsilon(Q)=I_\varepsilon(M),\qquad
 A_m(w_n)+\varepsilon A_k(w_n)\longrightarrow0.
\]
Thus convergence is strong in $H^{k/2}$.

A real minimizer can be obtained without a diamagnetic inequality.
The real and complex infima coincide: write $f=A+iB$ and, if both
component masses are positive, put
\[
 \vartheta=M(A)/M(f),\qquad
 g=A/\sqrt\vartheta,\qquad h=B/\sqrt{1-\vartheta}.
\]
Both $g,h$ are real and have mass $M(f)$. The two real-even kinetic
forms split, while convexity gives
\[
 |f|^p=\big(\vartheta g^2+(1-\vartheta)h^2\big)^{p/2}
             \le\vartheta|g|^p+(1-\vartheta)|h|^p.
\]
Therefore $E_\varepsilon(f)\ge
\vartheta E_\varepsilon(g)+(1-\vartheta)E_\varepsilon(h)$,
which is at least the real infimum. The cases of a zero component
are immediate. A real minimizing sequence and the preceding
compactness proof give a real minimizer $Q_M$.

\emph{Stationary identities.}
The energy is $C^1$ on $H^{k/2}$. The Lagrange multiplier theorem
gives~\eqref{higher:profile} weakly for a real $\omega_M$.
Minimizing along the mass-preserving dilation gives directly
\begin{equation}\label{higher:pohozaev}
 \frac m2A_m(Q_M)+\frac{k\varepsilon}2A_k(Q_M)
                         -\frac mp\|Q_M\|_p^p=0.
\end{equation}
Only the explicit scalar formula~\eqref{higher:dilation} is
differentiated; no spatial moment or energy-space dilation generator
is needed. Equation~\eqref{higher:energy} follows immediately.
Testing the weak equation against $Q_M$ additionally gives
\begin{equation}\label{higher:frequency}
 \omega_M M=\frac mN A_m(Q_M)
   +\left(\frac{k}{m}+\frac{k}{N}-1\right)
           \varepsilon A_k(Q_M)>0.
\end{equation}

\emph{Regularity.}
For $\omega_M>0$, the resolvent
\[
 R=(|D|^m+\varepsilon|D|^k+\omega_M)^{-1}
\]
maps $L^q$ to $H^{k,q}$ for every $1<q<\infty$. Its multiplier
after multiplication by $(1+|\xi|^2)^{k/2}$ satisfies the Mihlin
bounds, including at low frequency after each derivative is weighted
by the corresponding power of $|\xi|$. The equation reads
$Q_M=R(|Q_M|^{2m/N}Q_M)$.
If $N<k$, the starting energy embedding gives $L^\infty$; if $N=k$,
arbitrarily large finite input exponents give the same conclusion.
If $N>k$, start at
\[
 \beta_0=\frac12-\frac{k}{2N},\qquad
 \beta_{j+1}=(1+2m/N)\beta_j-k/N,
\]
where $\beta_j$ is the reciprocal integrability exponent. The fixed
point is $k/(2m)>1/2>\beta_0$, so finitely many iterations give an
input exponent above $N/k$ and hence $L^\infty$. The first nonlinear
input exponent exceeds one because $2N/(N-k)>1+2m/N$; all later
ones improve. An endpoint can be bypassed by slightly weakening an
intermediate exponent. Boundedness gives the nonlinearity in $L^2$,
so $Q_M\in H^k\subset H^1$. The $C^1$ chain rule then puts the
nonlinearity in $H^1$, and one more resolvent step gives
$Q_M\in H^{k+1}$. This proves the theorem.
\end{proof}

\subsection{Quantitative degeneration at the lower-order threshold}

\begin{proposition}[Low-frequency bounds]\label{higher:low-frequency}
For the minimizers in Theorem~\ref{higher:existence}, set
\[d=(M/M_{c,m})^{m/N}-1>0.\]
Then
\begin{align}
 \varepsilon A_k(Q_M)&\le\frac mk d A_m(Q_M),
       \label{higher:ratio}\\
 A_k(Q_M)&\le M\left(\frac{m d}{k\varepsilon}\right)^{k/(k-m)},
 & A_m(Q_M)&\le M\left(\frac{m d}{k\varepsilon}\right)^{m/(k-m)},
       \label{higher:kinetic-bounds}\\
 -\frac{k-m}{2k}M d
  \left(\frac{m d}{k\varepsilon}\right)^{m/(k-m)}
     &\le I_\varepsilon(M)<0.
       \label{higher:energy-bound}
\end{align}
Thus $A_m$, $A_k$, $|E_\varepsilon|$, and $\omega_M$ tend to zero
as $M\downarrow M_{c,m}$ with $\varepsilon$ fixed. After normalizing
$A_m=1$, these minimizers converge subsequentially in $H^{m/2}$,
modulo translations and constant phases, to a lower-order sharp
optimizer.
\end{proposition}
\begin{proof}
Combining~\eqref{higher:gn} and~\eqref{higher:pohozaev} gives
\eqref{higher:ratio}. Fourier interpolation gives
$A_m\le M^{1-m/k}A_k^{m/k}$, which proves
\eqref{higher:kinetic-bounds}; the energy identity gives
\eqref{higher:energy-bound}. Equation~\eqref{higher:frequency}
then gives $\omega_M\to0$.

For $M_n\downarrow M_{c,m}$, put
\[
 \lambda_n=A_m(Q_{M_n})^{-1/m},\qquad
 \widetilde Q_n(x)=\lambda_n^{N/2}Q_{M_n}(\lambda_nx).
\]
These profiles satisfy $A_m(\widetilde Q_n)=1$, $M(\widetilde Q_n)=M_n$,
and
\[
 E_0(\widetilde Q_n)
 =-\frac{k\varepsilon A_k(Q_{M_n})}{2mA_m(Q_{M_n})}<0,
 \qquad |E_0(\widetilde Q_n)|\le d_n/2\longrightarrow0.
\]
The normalized sharp-GN compactness lemma therefore gives the
asserted strong convergence. The higher-order coefficient in the
normalized equation is
\[
 \varepsilon\lambda_n^{m-k}
       =\varepsilon A_m(Q_{M_n})^{(k-m)/m}\longrightarrow0.
\]
Thus this term becomes small but still determines the negative
energy in~\eqref{higher:energy}.
\end{proof}

If a lower-order optimizer $Q$ additionally belongs to $H^{k/2}$,
one obtains matching energy powers by using the trial family
$\sqrt{M/M_{c,m}}\,Q_\lambda$. Put $a=A_m(Q)$, $b=A_k(Q)$,
and $\vartheta=M/M_{c,m}$, with $M(Q)=M_{c,m}$. Its optimal scale
is $\lambda^{k-m}=m a d/(k\varepsilon b)$, yielding
\begin{equation}\label{higher:energy-upper}
 I_\varepsilon(M)\le
 -\frac{\vartheta(k-m)}{2k}a d
       \left(\frac{m a d}{k\varepsilon b}\right)^{m/(k-m)}.
\end{equation}
Together with~\eqref{higher:energy-bound}, this gives the same power
$\varepsilon^{-m/(k-m)}d^{k/(k-m)}$ on both sides, not an exact
asymptotic coefficient. The additional optimizer regularity is
retained here for arbitrarily large $k$; the existence theorem does
not require it.

\begin{remark}[Meaning of the boundary]
For $b(\xi)=\varepsilon|\xi|^k$ one has
$\xi\cdot\nabla b=kb>mb$ away from zero. This is precisely the
opposite dilation sign and accounts for~\eqref{higher:energy}.
Every positive coefficient admits negative-energy stationary waves
at masses arbitrarily close to $M_{c,m}$. This modification also violates
the order-$m$ growth bound, so by itself it does not isolate the role of
the upper dilation condition within that growth class. A small
coefficient here is not a small bounded quadratic-form perturbation on
$H^{m/2}$, and the modified symbol is not homogeneous of degree $m$;
the fixed-order homogeneous robustness theorem is unaffected.

For a strong $H^{k/2}$ flow with conservation and the usual norm
continuation criterion,~\eqref{higher:coercivity} additionally rules
out finite-time divergence of the energy norm. This conditional
flow consequence is separate from the explicitly global stationary
solutions proved above.
\end{remark}


\begin{thebibliography}{9}
\bibitem{FJMM} A. J. Fern\'andez, L. Jeanjean, R. Mandel and M. Mari\c{s}, \emph{Non-homogeneous Gagliardo--Nirenberg inequalities in $\mathbb R^N$ and application to a biharmonic non-linear Schr\"odinger equation}, J. Differential Equations \textbf{330} (2022), 1--65, \href{https://arxiv.org/abs/2010.01448}{arXiv:2010.01448v3}.
\bibitem{IPS} V. Iyer, R. Parker and A. G. Stefanov, \emph{Existence and stability for traveling waves of fourth order semilinear wave and Schr\"odinger equations}, \href{https://arxiv.org/abs/2512.12390}{arXiv:2512.12390}, Section 3.3.2.
\bibitem{CW} W. Chen and Z. Wang, \emph{Normalized solutions to the biharmonic nonlinear Schr\"odinger equation with combined nonlinearities}, \href{https://arxiv.org/abs/2305.16565}{arXiv:2305.16565} (2023).
\bibitem{BCMN} D. Bonheure, J.-B. Cast\'eras, E. Moreira dos Santos and R. Nascimento, \emph{Orbitally stable standing waves of a mixed dispersion nonlinear Schr\"odinger equation}, SIAM J. Math. Anal. \textbf{50} (2018), 5027--5071, \href{https://arxiv.org/abs/1710.09775}{arXiv:1710.09775}.
\bibitem{LH} T. Luo and H. Hajaiej, \emph{Normalized solutions for a class of scalar field equations involving mixed fractional Laplacians}, Adv. Nonlinear Stud. \textbf{22} (2022), 228--247, \href{https://doi.org/10.1515/ans-2022-0013}{doi:10.1515/ans-2022-0013}.
\bibitem{PS} B. Pausader and S. Shao, \emph{The mass-critical fourth-order Schr\"odinger equation in high dimensions}, \href{https://arxiv.org/abs/0906.1547}{arXiv:0906.1547}, Section 6.
\bibitem{IKN} T. Inui, K. Komada and K. Nishimura, \emph{Scattering of non-radial solutions for the fourth-order Schr\"odinger equation with $\dot H^{1/2}$-critical and supercritical nonlinearity}, \href{https://arxiv.org/abs/2610.05802}{arXiv:2610.05802v1} (5 October 2026).
\end{thebibliography}
\end{document}
